\documentclass[12pt,a4paper]{article}
\usepackage[utf8]{inputenc}
\usepackage[polish]{babel}
\usepackage[T1]{fontenc}
\usepackage{geometry}
\usepackage{graphicx}
\usepackage{enumitem}
\usepackage{titlesec}
\usepackage{hyperref}
\usepackage{changepage}
\usepackage{float}
\usepackage{listings}
\usepackage{tcolorbox}
\usepackage{fontawesome5}
\usepackage{array}
\usepackage{longtable}
\usepackage{booktabs}
\usepackage{amsmath}
% Long code names (with _ and /) may break at line ends
\DeclareRobustCommand{\_}{\textunderscore\allowbreak}
\setlength{\emergencystretch}{2em}
\lstset{basicstyle=\ttfamily\small, columns=fullflexible, breaklines=true}

% UML diagrams drawn in LaTeX (TikZ, pgf-umlcd, pgf-umlsd)
\usepackage{tikz}
\usetikzlibrary{positioning,calc,shapes.geometric,arrows.meta,babel}
\usepackage[simplified]{pgf-umlcd}
\usepackage[underline=false]{pgf-umlsd}
\renewcommand{\umldrawcolor}{black}
\renewcommand{\umlfillcolor}{white}
\renewcommand{\umltextcolor}{black}
\tikzset{
    umlcd style class/.append style={align=left, inner sep=3pt},
    every picture/.append style={execute at begin picture={\hbadness=10000}},
    usecase/.style={ellipse, draw, minimum width=4.6cm, minimum height=1.05cm,
        align=center, inner sep=1.5pt, font=\footnotesize},
    sysbox/.style={draw, rectangle, inner sep=0pt},
    mult/.style={font=\scriptsize, inner sep=1.5pt},
    rolelabel/.style={font=\scriptsize\itshape, inner sep=1.5pt},
    dep/.style={dashed, -{Straight Barb[length=2mm]}},
    comp/.style={{Diamond[fill=black,length=3mm,width=2mm]}-},
}
% One class attribute per line, never broken
\newcommand{\attr}[1]{\attribute{\mbox{#1}}}
% Stick figure of a UML actor: \umlactor{name}{position}{label};
% lines start at the hands: (name r) and (name l)
\newcommand{\umlactor}[3]{%
    \begin{scope}[shift={(#2)}]
        \draw (0,0.62) circle (0.17);
        \draw (0,0.45) -- (0,-0.05) (-0.3,0.27) -- (0.3,0.27)
              (0,-0.05) -- (-0.22,-0.45) (0,-0.05) -- (0.22,-0.45);
        \node[below, font=\footnotesize, align=center] at (0,-0.45) {#3};
        \coordinate (#1 r) at (0.3,0.27);
        \coordinate (#1 l) at (-0.3,0.27);
    \end{scope}%
}

\geometry{
    a4paper,
    left=2.5cm,
    right=2.5cm,
    top=2.5cm,
    bottom=2.5cm
}

\titleformat{\section}{\normalfont\Large\bfseries}{\thesection}{1em}{}
\titleformat{\subsection}{\normalfont\large\bfseries}{\thesubsection}{1em}{}
\titleformat{\subsubsection}{\normalfont\normalsize\bfseries}{\thesubsubsection}{1em}{}

\begin{document}

\begin{titlepage}
\vspace{6cm}
    \centering

    \begin{figure}[h]
        \centering
        \includegraphics[width=0.3\linewidth]{images/uwbLogo.jpg}
        
    \end{figure}
    
    {\LARGE Uniwersytet w Białymstoku\\}
    
    \vspace{0.5cm}
    {\Large Wydział Informatyki\\}
    \vspace{0.3cm}
    {\large Kierunek: Informatyka II stopnia\\}
    {\large Rok I semestr I\\}
    \vspace{1.5cm}
    
    {\Large Inteligentne multimedialne systemy uczące\\}
    \vspace{1cm}
    
    {\huge\bfseries IMSU\\}
    
    \vspace{3cm}
    
    \begin{flushleft}
    {\large Prowadzący:\\}
    {\large dr inż. Wiesław Półjanowicz\\}
    \vspace{1cm}
    
    {\large Wykonali:\\}
    
    {\large Patryk Kulesza\\}
    {\large Bartłomiej Koźluk\\}
    {\large Mateusz Malewski\\}
    \end{flushleft}
    
    \vfill
    
    {\large Białystok, dn. \ldots\ldots\ldots\ldots\ldots}
\end{titlepage}

\tableofcontents
\newpage

\section{Treść zadania projektowego}

\textbf{IMSU -- multimedialny inteligentny system uczący:}

Klasyczne kursy internetowe podają materiał każdemu uczniowi w ten sam sposób, choć jedni lepiej uczą się z rysunków i schematów, inni z nagrań i rozmowy, a jeszcze inni przez samodzielne wykonywanie ćwiczeń. Nauczyciel prowadzący zajęcia na żywo dostosowuje się do ucznia odruchowo -- system internetowy musi to robić według jawnych reguł.

Celem projektu jest system, w którym \textbf{nauczyciel} przygotowuje kurs z lekcji, a każda lekcja ma tę samą treść w \textbf{trzech wersjach}: wzrokowej (V -- tekst ze schematami i wykresami), słuchowej (A -- nagranie z krótkim opisem) i kinestetycznej (K -- instrukcja krok po kroku z kartą zadania PDF), oraz jeden quiz z \textbf{puli} pytań z wyjaśnieniami. \textbf{Uczeń} po rejestracji wypełnia raz ankietę VAK, z której system wyznacza jeden aktywny styl i profil procentowy V / A / K. Uczeń uczy się zawsze w jednej formie naraz: dostaje wersję lekcji w formie swojego stylu, a może przerobić lekcję w innej formie. Lekcje są dostępne po kolei -- następna po zaliczeniu quizu poprzedniej.

Po każdym quizie system zmienia procenty form według jawnych reguł, a gdy uczniowi nie idzie, proponuje przerobienie lekcji w innej formie albo zmianę stylu na stałe -- zawsze z liczbami, z których ta decyzja wynika. Po zmianie stylu okres próbny sprawdza, czy zmiana pomogła. Każda decyzja trafia do dziennika, który widzi uczeń i nauczyciel kursu; o zmianie stylu decyduje uczeń.

System obsługuje trzy role: ucznia, nauczyciela (konto nauczyciela zatwierdza administrator) i administratora, który zarządza kontami użytkowników. Uczeń prosi o zapis na kurs, a autor kursu prośbę przyjmuje albo odrzuca.

\section{Cel budowania systemu i zakres systemu}

\subsection{Cele główne}

\begin{itemize}[leftmargin=*]
    \item Nauka w jednym stylu VAK naraz: każda lekcja w trzech wersjach (V / A / K), uczeń dostaje wersję w formie swojego stylu
    \item Sprawdzanie stylu wynikami: ankieta jest hipotezą startową, a procenty form zmieniają wyniki quizów
    \item Reakcja na trudności z argumentami z danych: sugestia przerobienia lekcji w innej formie, propozycja zmiany stylu na stałe, okres próbny
    \item Wyjaśnialność: każda reakcja systemu wynika z reguły w czystej funkcji, niesie liczby i trafia do dziennika decyzji
\end{itemize}

\subsection{Cele szczegółowe}

\begin{itemize}[leftmargin=*]
    \item Profil procentowy z ankiety (suma 100, każda forma co najmniej 5\%), ankieta wypełniana raz; przy remisie uczeń sam wybiera jeden styl
    \item Zmiana procentów tylko po pierwszym podejściu w danej formie w danej lekcji, żeby wyników nie dało się ,,nabić'' powtórkami
    \item Propozycja zmiany stylu tylko dla formy \textbf{potwierdzonej wynikami} (co najmniej jedno zaliczenie w tej formie); decyzję podejmuje uczeń
    \item Quiz losowany z puli (5 pytań z co najmniej 5), oceniany wyłącznie na serwerze, z wyjaśnieniem przy każdym pytaniu
    \item Kolejność lekcji egzekwowana na serwerze: lekcja $n+1$ po zaliczeniu lekcji $n$
    \item Publikacja lekcji tylko wtedy, gdy ma zarys, wnioski, trzy kompletne wersje i co najmniej 5 pytań w puli
    \item Statystyki: ,,Moje postępy'' ucznia oraz panel kursu nauczyciela ze skutecznością form
    \item Prosty i czytelny interfejs, przystosowany do młodszych użytkowników
\end{itemize}

\subsection{Zakres systemu}

\begin{itemize}[leftmargin=*]
    \item Rejestracja z wyborem roli (uczeń lub nauczyciel), logowanie, sesja; zatwierdzanie kont nauczycieli przez administratora
    \item Administracja kont nauczycieli i uczniów (tworzenie, edycja, blokada, zatwierdzanie i odrzucanie)
    \item Kursy: wstęp, podsumowanie i wnioski, lekcje w ustalonej kolejności, publikacja, archiwizacja; zapisy na kurs z akceptacją autora
    \item Lekcja: zarys, wersje V / A / K (bloki tekstu i załączniki: obraz, film, nagranie, PDF, link), wnioski; gotowość i publikacja
    \item Quiz lekcji: pula 1--50 pytań jednokrotnego wyboru z wyjaśnieniami, losowanie 5 pytań na podejście
    \item Model ucznia: ankieta VAK, styl, profil procentowy i jego historia, wyniki quizów, wybory formy lekcji, propozycje zmiany stylu, okresy próbne, dziennik decyzji
    \item Statystyki ucznia (wykres procentów, postęp kursów, wyniki lekcji) i kursu (uczniowie, skuteczność form, rozkład stylów, szczegóły ucznia)
\end{itemize}

Poza zakresem: ponowna ankieta, test końcowy kursu, pytania inne niż jednokrotnego wyboru, powiadomienia e-mail, moderacja kursów przez administratora, modele językowe.

\section{Perspektywa wariantów użycia systemu}

\subsection{Diagram wariantów użycia systemu}

\begin{figure}[H]
    \centering
    % Use case diagram (plain TikZ)
\begin{tikzpicture}[font=\footnotesize]
    % System boundary
    \draw[sysbox] (2.2,0.95) rectangle (13.1,-11.65);
    \node[anchor=north, font=\small\bfseries] at (7.65,0.9) {System IMSU};

    % Student side
    \node[usecase] (uc1) at (4.75,-0.35) {UC-01 Rejestracja};
    \node[usecase] (uc2) at (4.75,-1.65) {UC-02 Logowanie};
    \node[usecase] (uc3) at (4.75,-2.95) {UC-03 Ankieta VAK\\(raz, remis $\to$ wybór)};
    \node[usecase] (uc4) at (4.75,-4.25) {UC-04 Katalog\\i prośba o zapis};
    \node[usecase] (uc5) at (4.75,-5.55) {UC-05 Nauka\\z lekcji};
    \node[usecase] (uc6) at (4.75,-6.85) {UC-06 Quiz z puli};
    \node[usecase, fill=black!6] (ad) at (4.75,-8.4) {Adaptacja profilu\\(reguły R2--R7)};
    \node[usecase] (uc7) at (4.75,-9.7) {UC-07 Zarządzanie\\stylem};
    \node[usecase] (uc8) at (4.75,-10.95) {UC-08 Moje postępy};

    % Teacher and administrator side
    \node[usecase] (uc9)  at (10.55,-0.95) {UC-09 Edycja kursu};
    \node[usecase] (uc10) at (10.55,-2.25) {UC-10 Edycja lekcji\\(3 wersje, pula)};
    \node[usecase] (uc11) at (10.55,-3.55) {UC-11 Przyjmowanie\\próśb o zapis};
    \node[usecase] (uc12) at (10.55,-4.85) {UC-12 Panel kursu\\i szczegóły ucznia};
    \node[usecase] (uc13) at (10.55,-8.4) {UC-13 Zatwierdzanie\\nauczycieli};
    \node[usecase] (uc14) at (10.55,-9.7) {UC-14 Zarządzanie\\kontami};
    \node[usecase] (uc15) at (10.55,-11.0) {UC-15 Zmiana\\właściciela kursu};

    % Actors
    \umlactor{guest}{0.75,-1.2}{Gość}
    \umlactor{student}{0.75,-6.6}{Uczeń}
    \umlactor{teacher}{14.45,-3.1}{Nauczyciel}
    \umlactor{admin}{14.45,-9.7}{Administrator}

    % Associations
    \foreach \uc in {uc1,uc2} \draw (guest r) -- (\uc.west);
    \foreach \uc in {uc3,uc4,uc5,uc6,uc7,uc8} \draw (student r) -- (\uc.west);
    \foreach \uc in {uc9,uc10,uc11,uc12} \draw (teacher l) -- (\uc.east);
    \foreach \uc in {uc13,uc14,uc15} \draw (admin l) -- (\uc.east);

    % Include
    \draw[dep] (uc6.south) -- node[right, font=\scriptsize] {«include»} (ad.north);
\end{tikzpicture}

    \caption{Diagram przypadków użycia systemu}
\end{figure}

Wariant ,,Adaptacja profilu'' nie ma własnego aktora: wykonuje go serwer przy każdym wysłaniu quizu (reguły R2--R7, rozdział~\ref{sec:rules}).

\subsection{Opisy tekstowe aktorów}

\begin{itemize}[leftmargin=*]
    \item \textbf{Gość} -- osoba niezalogowana. Widzi stronę startową z opisem systemu, może założyć konto (jako uczeń albo nauczyciel) lub zalogować się.

    \item \textbf{Uczeń} -- zalogowany użytkownik, który się uczy. Przed rozpoczęciem nauki wypełnia raz ankietę VAK (przy remisie wybiera jeden styl). Prosi o zapis na kursy z katalogu, uczy się z lekcji w formie swojego stylu, rozwiązuje quizy, decyduje o propozycjach zmiany stylu i ogląda swoje postępy oraz dziennik decyzji.

    \item \textbf{Nauczyciel} -- tworzy kursy i lekcje (trzy wersje, pula pytań), publikuje je, przyjmuje prośby o zapis i śledzi wyniki uczniów w panelu kursu. Konto nauczyciela założone przez rejestrację \textbf{wymaga zatwierdzenia} przez administratora; do tego czasu nauczyciel widzi tylko ekran ,,Konto czeka na zatwierdzenie'', a serwer odrzuca jego żądania kodem \texttt{403 TEACHER\_NOT\_APPROVED}.

    \item \textbf{Administrator} -- zatwierdza albo odrzuca konta nauczycieli oraz zarządza kontami nauczycieli i uczniów: tworzy je, edytuje i blokuje. Konto administratora tworzy się poleceniem w konsoli, nie przez rejestrację.
\end{itemize}

\subsection{Opisy tekstowe wariantów użycia}

\subsubsection{Rejestracja}
\textbf{Numer:} UC-01 \quad \textbf{Aktor:} Gość

\textbf{Krótki opis:} Założenie konta ucznia albo nauczyciela. Rola jest wybierana przy rejestracji i nie zmienia się później.

\textbf{Główny przepływ zdarzeń:}
\begin{enumerate}[leftmargin=*]
    \item Gość klika ,,Załóż konto'' na stronie startowej
    \item Gość wybiera rolę (Uczeń / Nauczyciel), podaje imię, e-mail i hasło (min. 8 znaków)
    \item System sprawdza dane i unikalność adresu e-mail
    \item System tworzy konto i rozpoczyna sesję; konto ucznia jest od razu aktywne, konto nauczyciela dostaje stan \texttt{pending} (czeka na zatwierdzenie, UC-13)
\end{enumerate}

\subsubsection{Logowanie}
\textbf{Numer:} UC-02 \quad \textbf{Aktor:} Gość

\textbf{Krótki opis:} Zalogowanie się adresem e-mail i hasłem.

\textbf{Główny przepływ zdarzeń:}
\begin{enumerate}[leftmargin=*]
    \item Gość klika ,,Zaloguj się'' i podaje e-mail oraz hasło
    \item System weryfikuje dane oraz to, czy konto nie jest zablokowane
    \item System zapisuje sesję w ciasteczku i przekierowuje do panelu użytkownika; uczeń bez ankiety trafia na ankietę, nauczyciel niezatwierdzony -- na ekran oczekiwania (albo odrzucenia)
\end{enumerate}

\subsubsection{Wypełnienie ankiety VAK}
\textbf{Numer:} UC-03 \quad \textbf{Aktor:} Uczeń

\textbf{Krótki opis:} Ankieta wyznacza jeden styl ucznia i startowy profil procentowy V / A / K (reguła R1). Ankietę wypełnia się raz -- potem procenty zmieniają tylko wyniki quizów.

\textbf{Główny przepływ zdarzeń:}
\begin{enumerate}[leftmargin=*]
    \item System przekierowuje ucznia bez ankiety na stronę ankiety (kursy i lekcje zwracają \texttt{403 SURVEY\_REQUIRED})
    \item Uczeń odpowiada na 20 pytań (każda odpowiedź odpowiada jednej formie; klient nie zna przypisania)
    \item Serwer zlicza punkty, zamienia je na procenty (np. A 11 / V 6 / K 3 $\to$ A 55 / V 30 / K 15) i wyznacza styl -- formę z największą liczbą punktów
    \item System pokazuje procenty na paskach i styl
\end{enumerate}

\textbf{Przepływ alternatywny -- remis:} kilka form ma tyle samo punktów; profil zostaje zapisany bez stylu, a uczeń wybiera jeden z remisujących stylów (,,Wybieram ten styl''). Do tego wyboru kursy są niedostępne. Ponowne wysłanie ankiety zwraca \texttt{409 SURVEY\_ALREADY\_DONE}.

\subsubsection{Katalog kursów i prośba o zapis}
\textbf{Numer:} UC-04 \quad \textbf{Aktor:} Uczeń

\textbf{Krótki opis:} Uczeń przegląda katalog kursów opublikowanych i prosi o zapis; treść kursu jest dostępna po przyjęciu prośby przez nauczyciela.

\textbf{Główny przepływ zdarzeń:}
\begin{enumerate}[leftmargin=*]
    \item Uczeń wybiera ,,Kursy'' -- system wyświetla sekcję ,,Moje kursy'' (przyjęte, z postępem ,,2 z 4 lekcji'') i ,,Katalog'' (tytuł, opis, nauczyciel, liczba lekcji)
    \item Uczeń klika ,,Zapisz się'' -- system zapisuje prośbę (\texttt{pending}) i pokazuje ,,Prośba wysłana -- czeka na nauczyciela''
    \item Po przyjęciu prośby (UC-11) kurs przechodzi do ,,Moich kursów''; strona kursu pokazuje wstęp, ponumerowane lekcje z zarysami i stanem (zaliczona, dostępna, zablokowana) oraz podsumowanie
\end{enumerate}

\textbf{Przepływ alternatywny:} prośba odrzucona -- system pokazuje ,,Prośba odrzucona'' i przycisk ,,Poproś ponownie''. Uczeń przyjęty, który nie ma jeszcze wyników w kursie, widzi podpowiedź z wyników innych uczniów (reguła R8), np. ,,W tym kursie najlepiej wypada forma słuchowa (średnio 86\%, 3 uczniów)''.

\subsubsection{Nauka z lekcji}
\textbf{Numer:} UC-05 \quad \textbf{Aktor:} Uczeń

\textbf{Krótki opis:} Uczeń otwiera lekcję w wersji odpowiadającej jego stylowi i w razie potrzeby przerabia ją w innej formie.

\textbf{Główny przepływ zdarzeń:}
\begin{enumerate}[leftmargin=*]
    \item Uczeń otwiera dostępną lekcję kursu, na który jest przyjęty
    \item System wyświetla zarys, pasek ,,Przerabiasz tę lekcję w formie: Wzrokowej (V)'', wersję lekcji w tej formie, wnioski i sekcję ,,Sprawdź się''
    \item Uczeń może wybrać ,,Przerób tę lekcję w innej formie'' -- system zapisuje wybór formy tylko dla tej lekcji, pokazuje inną wersję i zapisuje zdarzenie w dzienniku; styl ucznia się nie zmienia
\end{enumerate}

\textbf{Przepływ alternatywny -- blokada kolejności:} lekcja $n+1$ jest dostępna dopiero po zaliczeniu quizu lekcji $n$ (reguła R10). Otwarcie zablokowanej lekcji kończy się kodem \texttt{403 LESSON\_LOCKED} z numerem wymaganej lekcji, a klient pokazuje kłódkę i link ,,Najpierw zalicz lekcję $n$''.

\subsubsection{Rozwiązanie quizu z puli}
\textbf{Numer:} UC-06 \quad \textbf{Aktor:} Uczeń

\textbf{Krótki opis:} Uczeń rozwiązuje 5 pytań wylosowanych z puli lekcji; serwer ocenia odpowiedzi i aktualizuje model ucznia ($<<$include$>>$ Adaptacja profilu).

\textbf{Główny przepływ zdarzeń:}
\begin{enumerate}[leftmargin=*]
    \item Serwer losuje 5 pytań z puli (przy powtórce najpierw pytania spoza poprzedniego zestawu -- R9) i zapisuje zestaw; klient dostaje pytania bez klucza
    \item Uczeń zaznacza odpowiedzi i klika ,,Sprawdź''
    \item Serwer ocenia odpowiedzi wobec zapisanego zestawu; wynik co najmniej 70\% (4 z 5) oznacza zaliczenie
    \item Serwer wykonuje kroki adaptacji (R2--R7): zmiana procentów, odblokowanie następnej lekcji, pochwała, sugestia, okres próbny, propozycja
    \item System wyświetla wynik, przy każdym pytaniu wybraną i poprawną odpowiedź z \textbf{wyjaśnieniem} oraz ,,Co to zmieniło'', np. ,,Słuchowe 47\% $\to$ 57\% · Lekcja 4 odblokowana''
\end{enumerate}

\textbf{Przepływy alternatywne:}
\begin{itemize}[leftmargin=*]
    \item Quiz niezaliczony -- ramka sugestii z uzasadnieniem i przyciskiem ,,Przerób tę lekcję w formie X'' (R5) oraz ,,Rozwiąż jeszcze raz'' (nowe losowanie)
    \item Kolejne podejście w tej samej formie -- ,,To kolejne podejście w tej formie -- procenty bez zmian''; zaliczenie odblokowuje jednak następną lekcję
    \item Forma potwierdzona wynikami ma wyższy procent niż styl -- baner propozycji zmiany stylu (UC-07)
\end{itemize}

\subsubsection{Zarządzanie stylem}
\textbf{Numer:} UC-07 \quad \textbf{Aktor:} Uczeń

\textbf{Krótki opis:} Uczeń przegląda swój styl i profil, odpowiada na propozycje zmiany stylu i może zmienić styl ręcznie.

\textbf{Główny przepływ zdarzeń:}
\begin{enumerate}[leftmargin=*]
    \item Uczeń wybiera ,,Mój styl'' -- system wyświetla aktywny styl i jego źródło, procenty V / A / K z opisem przy każdej formie (,,potwierdzona wynikami (2 lekcje, średnio 88\%)'' albo ,,tylko z ankiety'') i wynik ankiety
    \item Jeśli jest otwarta propozycja, system pokazuje ją z dowodem, np. ,,Wzrokowe: 3 lekcje, średnio 53\% · Słuchowe: 1 lekcja, średnio 80\%''
    \item Uczeń klika ,,Przełącz na stałe'' -- styl się zmienia, wybory formy w niezaliczonych lekcjach są czyszczone, zaczyna się \textbf{okres próbny} 2 lekcji (R7); albo ,,Zostaję przy swoim'' -- propozycja tej samej formy wraca najwcześniej po 2 nowych policzonych wynikach
    \item Po okresie próbnym system pokazuje jego wynik, np. ,,Zmiana pomogła: średnia 40\% $\to$ 90\%''
\end{enumerate}

\textbf{Przepływ alternatywny -- ręczna zmiana:} uczeń wybiera inny styl i potwierdza w oknie (,,Procenty się nie zmienią, zmiana trafi do dziennika''). Otwarta propozycja dostaje stan \texttt{superseded}, zaczyna się okres próbny.

\subsubsection{Moje postępy}
\textbf{Numer:} UC-08 \quad \textbf{Aktor:} Uczeń

\textbf{Krótki opis:} Uczeń ogląda historię swojego profilu, postęp kursów, wyniki lekcji i dziennik decyzji.

\textbf{Główny przepływ zdarzeń:}
\begin{enumerate}[leftmargin=*]
    \item Uczeń wybiera ,,Moje postępy''
    \item System wyświetla wykres procentów V / A / K w czasie (z tabelą tych samych danych), postęp kursów (,,następna: Lekcja 3'') i tabelę wyników lekcji z formą każdego podejścia (pierwsze policzone wyróżnione)
    \item System wyświetla dziennik decyzji od najnowszej: data, zdanie złożone z kodu i liczb zdarzenia, kurs i lekcja
\end{enumerate}

\subsubsection{Edycja kursu}
\textbf{Numer:} UC-09 \quad \textbf{Aktor:} Nauczyciel

\textbf{Krótki opis:} Nauczyciel tworzy i edytuje swoje kursy oraz decyduje o ich publikacji.

\textbf{Główny przepływ zdarzeń:}
\begin{enumerate}[leftmargin=*]
    \item Nauczyciel wybiera ,,Moje kursy'' i klika ,,Nowy kurs''; podaje tytuł i opis -- system tworzy kurs prywatny
    \item Nauczyciel uzupełnia wstęp oraz ,,Podsumowanie i wnioski'', dodaje lekcje i ustala ich kolejność
    \item Nauczyciel publikuje kurs -- kurs pojawia się w katalogu uczniów; kursu bez opublikowanej lekcji nie da się opublikować (\texttt{409 COURSE\_HAS\_NO\_PUBLISHED\_LESSONS})
\end{enumerate}

\subsubsection{Edycja lekcji}
\textbf{Numer:} UC-10 \quad \textbf{Aktor:} Nauczyciel

\textbf{Krótki opis:} Nauczyciel układa lekcję w trzech wersjach, tworzy pulę pytań i publikuje lekcję, gdy jest kompletna.

\textbf{Główny przepływ zdarzeń:}
\begin{enumerate}[leftmargin=*]
    \item Nauczyciel otwiera lekcję w edytorze z zakładkami ,,Zarys i wnioski'', ,,Wzrokowa (V)'', ,,Słuchowa (A)'', ,,Kinestetyczna (K)'' i ,,Quiz (pula)''
    \item W zakładce wersji dodaje bloki tekstu i załączniki (obraz, film, nagranie, PDF, link); edytor podpowiada, co wersja musi zawierać
    \item W zakładce quizu dodaje pytania (3--4 odpowiedzi, jedna poprawna, wymagane wyjaśnienie); licznik pokazuje ,,8 pytań w puli (min. 5)''
    \item Panel gotowości pokazuje listę braków (R11); gdy lekcja jest kompletna, nauczyciel klika ,,Opublikuj lekcję''
\end{enumerate}

\textbf{Przepływy alternatywne:} publikacja niekompletnej lekcji -- \texttt{409 LESSON\_NOT\_READY} z listą braków; zapis, który zrobiłby opublikowaną lekcję niekompletną -- \texttt{409 PUBLISHED\_LESSON\_INCOMPLETE}; wycofanie ostatniej opublikowanej lekcji opublikowanego kursu -- \texttt{409 LAST\_PUBLISHED\_LESSON}.

\subsubsection{Przyjmowanie próśb o zapis}
\textbf{Numer:} UC-11 \quad \textbf{Aktor:} Nauczyciel

\textbf{Krótki opis:} Autor kursu decyduje, którzy uczniowie mają dostęp do treści kursu.

\textbf{Główny przepływ zdarzeń:}
\begin{enumerate}[leftmargin=*]
    \item Nauczyciel widzi przy kursie liczbę próśb (,,2 prośby o zapis'') i otwiera ,,Panel kursu''
    \item Sekcja ,,Prośby o zapis'' pokazuje imię, e-mail i datę prośby
    \item Nauczyciel klika ,,Przyjmij'' albo ,,Odrzuć''; przyjęty uczeń zachowuje dostęp także po wycofaniu publikacji kursu
\end{enumerate}

\subsubsection{Panel kursu i szczegóły ucznia}
\textbf{Numer:} UC-12 \quad \textbf{Aktor:} Nauczyciel

\textbf{Krótki opis:} Nauczyciel śledzi wyniki uczniów swojego kursu i uzasadnienia decyzji systemu.

\textbf{Główny przepływ zdarzeń:}
\begin{enumerate}[leftmargin=*]
    \item Nauczyciel otwiera ,,Panel kursu''
    \item System wyświetla tabelę uczniów (styl, procenty, postęp, średnia, ostatnia aktywność, otwarta propozycja), \textbf{skuteczność form} (dla V / A / K: średnia policzonych wyników, liczba wyników, podejść i uczniów) oraz rozkład stylów
    \item Nauczyciel klika ucznia -- system pokazuje jego profil, wykres procentów, wyniki lekcji tego kursu i dziennik decyzji
\end{enumerate}

\subsubsection{Zatwierdzanie nauczycieli}
\textbf{Numer:} UC-13 \quad \textbf{Aktor:} Administrator

\textbf{Krótki opis:} Administrator decyduje, czy konto nauczyciela z rejestracji może tworzyć kursy.

\textbf{Główny przepływ zdarzeń:}
\begin{enumerate}[leftmargin=*]
    \item Administrator wybiera ,,Użytkownicy'', zakładkę ,,Nauczyciele'' i filtr ,,Czekające (n)''
    \item Kolumna ,,Konto'' pokazuje stan tekstem: Zatwierdzone / Czeka na zatwierdzenie / Odrzucone
    \item Administrator klika ,,Zatwierdź'' albo ,,Odrzuć'' (odrzucenie z potwierdzeniem); zatwierdzony nauczyciel widzi ,,Moje kursy''
\end{enumerate}

\subsubsection{Zarządzanie kontami}
\textbf{Numer:} UC-14 \quad \textbf{Aktor:} Administrator

\textbf{Krótki opis:} Administrator zarządza kontami nauczycieli i uczniów.

\textbf{Główny przepływ zdarzeń:}
\begin{enumerate}[leftmargin=*]
    \item Administrator wybiera ,,Użytkownicy'' i zakładkę Nauczyciele lub Uczniowie
    \item Administrator dodaje konto (imię, e-mail, hasło) lub edytuje istniejące; nauczyciel założony przez administratora jest od razu zatwierdzony
    \item Administrator może zablokować konto -- zablokowany użytkownik nie zaloguje się, a jego sesja przestaje działać przy następnym żądaniu
\end{enumerate}

\subsection{Przykładowe diagramy sekwencji}

Przykłady 2 i 3 to ścieżka ucznia demonstracyjnego Kuby (\texttt{uczen.demo@imsu.pl}) z kursu ,,Matematyka: funkcja liniowa'', lekcja 3 (rozdział~\ref{sec:rules-example}).

\subsubsection{Przykład 1. Pierwsze logowanie ucznia i ankieta VAK}

\begin{figure}[H]
    \centering
    % Sequence diagram 1: first login and the VAK survey (R1)
\begin{sequencediagram}
    \scriptsize
    \def\unitfactor{0.72}
    \newthread{u}{Uczeń}
    \newinst[1.7]{k}{Klient (Next.js)}
    \newinst[2.1]{api}{Serwer (FastAPI)}
    \newinst[1.3]{r}{\texttt{rules.py}}
    \newinst[0.8]{db}{Baza danych}

    \begin{call}{u}{e-mail, hasło}{k}{ankieta}
        \begin{call}{k}{\texttt{POST /auth/login}}{api}{200 + ciasteczko}
        \end{call}
        \begin{call}{k}{\texttt{GET /profile/learning-style}}{api}{bez ankiety}
        \end{call}
        \begin{callself}{k}{przejście na \texttt{/survey}}{}
        \end{callself}
        \begin{call}{k}{\texttt{GET /survey/questions}}{api}{20 pytań bez klucza}
        \end{call}
    \end{call}

    \begin{call}{u}{20 odpowiedzi}{k}{A 55 / V 30 / K 15}
        \begin{call}{k}{\texttt{POST /survey/answers}}{api}{profil ucznia}
            \begin{call}{api}{A 11 / V 6 / K 3}{r}{A 55 / V 30 / K 15}
            \end{call}
            \begin{call}{api}{wynik, profil, historia}{db}{commit}
            \end{call}
        \end{call}
    \end{call}

    \begin{sdblock}{alt}{remis punktów, np. V 8 / A 8 / K 4}
        \begin{call}{u}{wybór stylu}{k}{dostęp do kursów}
            \begin{call}{k}{\texttt{POST .../tie-choice}}{api}{profil ze stylem}
            \end{call}
        \end{call}
    \end{sdblock}
\end{sequencediagram}

    \caption{Pierwsze logowanie i ankieta z procentami}
\end{figure}

Uczeń bez ankiety loguje się do systemu. Klient pobiera profil, a ponieważ ankieta nie została wypełniona (\texttt{survey\_completed = false}), przekierowuje ucznia na stronę ankiety. Pytania przychodzą bez informacji, której formie odpowiada dana odpowiedź -- punkty liczy serwer. Funkcja \texttt{survey\_percentages} z \texttt{rules.py} zamienia punkty A 11 / V 6 / K 3 na A 55 / V 30 / K 15, a \texttt{survey\_leaders} wskazuje jedną formę z maksimum punktów, więc styl to A. Serwer zapisuje w jednej transakcji wynik ankiety, profil, pierwszy punkt historii procentów i zdarzenie \texttt{survey\_completed}. Przy remisie profil nie ma stylu, a uczeń wybiera jeden z remisujących stylów (\texttt{POST /profile/learning-style/tie-choice}).

\subsubsection{Przykład 2. Niezaliczony quiz, sugestia i lekcja w innej formie}

\begin{figure}[H]
    \centering
    % Sequence diagram 2: failed quiz -> suggestion -> the lesson in another form (R2, R5, R9)
\begin{sequencediagram}
    \scriptsize
    \def\unitfactor{0.72}
    \newthread{k}{Uczeń (klient)}
    \newinst[2.6]{as}{\texttt{assessment}}
    \newinst[1.9]{ls}{\texttt{learning\_style}}
    \newinst[3.1]{db}{Baza danych}

    \begin{call}{k}{\texttt{GET /lessons/\{id\}/quiz}}{as}{5 z 8 pytań}
    \end{call}

    \begin{call}{k}{\texttt{POST .../quiz/attempts}}{as}{40\%, wyjaśnienia}
        \begin{callself}{as}{ocena: 2 z 5}{}
        \end{callself}
        \begin{call}{as}{\texttt{record\_quiz\_result}}{ls}{\texttt{QuizFeedback}}
            \begin{callself}{ls}{R2: V 45 $\to$ 35, A 47, K 18}{}
            \end{callself}
            \begin{callself}{ls}{R5: sugestia A (ankieta 40\%)}{}
            \end{callself}
            \begin{call}{ls}{wynik, profil, zdarzenia}{db}{commit}
            \end{call}
        \end{call}
    \end{call}

    \begin{call}{k}{\texttt{PUT .../study/form} (A)}{ls}{wersja A lekcji}
        \begin{call}{ls}{wybór formy, zdarzenie}{db}{commit}
        \end{call}
    \end{call}

    \begin{call}{k}{\texttt{GET /lessons/\{id\}/quiz}}{as}{nowy zestaw (R9)}
    \end{call}
\end{sequencediagram}

    \caption{Niezaliczony quiz $\to$ sugestia $\to$ przerobienie lekcji w innej formie}
\end{figure}

Kuba (styl V, profil V 45 / A 40 / K 15) dostaje 5 pytań wylosowanych z puli 8 i odpowiada poprawnie na 2. Moduł \texttt{assessment} ocenia zestaw zapisany na serwerze (40\%) i przekazuje wynik do \texttt{learning\_style.hooks.record\_quiz\_result}. To pierwsze podejście w formie V w tej lekcji, więc wynik jest policzony: V spada o 10 punktów, a A i K dzielą resztę proporcjonalnie -- V 35 / A 47 / K 18 (R2). Kuba nie ma zaliczonej formy innej niż V, więc sugestia (R5) wskazuje drugą formę z ankiety: A (40\%). Klient pokazuje wynik z wyjaśnieniami i ramkę sugestii. Po kliknięciu ,,Przerób tę lekcję w formie słuchowej'' serwer zapisuje wybór formy lekcji (\texttt{lesson\_forms}) i zdarzenie \texttt{lesson\_form\_changed}, a następne losowanie daje najpierw pytania spoza poprzedniego zestawu (R9).

\subsubsection{Przykład 3. Zaliczony quiz, propozycja zmiany stylu i okres próbny}

\begin{figure}[H]
    \centering
    % Sequence diagram 3: passed quiz -> style proposal -> trial period (R2, R4, R6, R7, R10)
\begin{sequencediagram}
    \scriptsize
    \def\unitfactor{0.72}
    \newthread{k}{Uczeń (klient)}
    \newinst[2.6]{as}{\texttt{assessment}}
    \newinst[1.7]{ls}{\texttt{learning\_style}}
    \newinst[2.9]{cu}{\texttt{curriculum}}

    \begin{call}{k}{\texttt{POST .../quiz/attempts}}{as}{80\%, propozycja}
        \begin{call}{as}{\texttt{record\_quiz\_result}}{ls}{\texttt{QuizFeedback}}
            \begin{callself}{ls}{R2: A 47 $\to$ 57, V 28, K 15}{}
            \end{callself}
            \begin{call}{ls}{\texttt{mark\_lesson\_passed}}{cu}{lekcja 4 (R10)}
            \end{call}
            \begin{callself}{ls}{R6: propozycja V $\to$ A}{}
            \end{callself}
            \begin{callself}{ls}{zapis, jeden commit}{}
            \end{callself}
        \end{call}
    \end{call}

    \begin{call}{k}{\texttt{POST .../proposals/\{id\}/accept}}{ls}{styl A, okres 0 z 2}
        \begin{callself}{ls}{R6: czyszczenie wyborów}{}
        \end{callself}
        \begin{callself}{ls}{R7: start, przed = 40\%}{}
        \end{callself}
    \end{call}

    \begin{sdblock}{loop}{2 kolejne policzone wyniki w formie A}
        \begin{call}{k}{\texttt{POST .../quiz/attempts}}{as}{wynik okresu}
            \begin{call}{as}{\texttt{record\_quiz\_result}}{ls}{}
                \begin{callself}{ls}{R7: po $-$ przed}{}
                \end{callself}
            \end{call}
        \end{call}
    \end{sdblock}
\end{sequencediagram}

    \caption{Zaliczony quiz $\to$ propozycja zmiany stylu $\to$ okres próbny}
\end{figure}

Kuba rozwiązuje quiz lekcji 3 w formie A na 80\%. Pierwsze podejście w formie A w tej lekcji jest policzone: A rośnie o 10 punktów -- V 28 / A 57 / K 15 (R2). Zaliczenie zapisuje \texttt{curriculum.mark\_lesson\_passed} i odblokowuje lekcję 4 (R10). Forma A ma teraz zaliczenie (jest potwierdzona wynikami, R3) i wyższy procent niż styl V (57 wobec 28), więc powstaje propozycja zmiany stylu V $\to$ A z dowodem ,,Wzrokowe: 3 lekcje, średnio 53\% · Słuchowe: 1 lekcja, średnio 80\%'' (R6). Wszystkie zapisy idą jednym \texttt{commit}. Po przyjęciu propozycji styl zmienia się na A, wybory formy w niezaliczonych lekcjach są czyszczone (u Kuby lista jest pusta -- lekcję 3 już zaliczył) i zaczyna się okres próbny ze średnią ,,przed'' 40\% (dwa ostatnie policzone wyniki w V: 40\% i 40\%). Po 2 policzonych wynikach w formie A serwer porównuje średnie (R7).

\section{Perspektywa logiczna}

\subsection{Diagram klas}

Diagram klas podzielono na trzy rysunki według modułów serwera. Klasy odpowiadają modelom SQLAlchemy z plików \nolinkurl{backend/app/modules/*/models.py}; typ zakończony znakiem \texttt{?} oznacza wartość opcjonalną (\texttt{null}). Klasa \texttt{User} bez atrybutów na rysunkach~\ref{fig:classes-content} i~\ref{fig:classes-student} to ta sama klasa co na rysunku~\ref{fig:classes-courses}, a \texttt{Lesson} na rysunku~\ref{fig:classes-content} -- ta sama co na rysunku~\ref{fig:classes-courses}. Wszystkie klasy modelu ucznia (rysunek~\ref{fig:classes-student}) odwołują się do ucznia przez \texttt{user\_id}; odwołania do kursu, lekcji i zestawu pytań są zapisane jako atrybuty \texttt{course\_id}, \texttt{lesson\_id} i \texttt{draw\_id}.

\begin{figure}[H]
    \centering
    % Class diagram 1: accounts, courses, enrollments (identity, curriculum)
\begin{tikzpicture}[font=\scriptsize]
    \begin{class}[text width=4.5cm]{User}{0,0}
        \attr{id : int}
        \attr{email : str}
        \attr{password\_hash : str}
        \attr{display\_name : str}
        \attr{role : UserRole}
        \attr{approval\_status : ApprovalStatus}
        \attr{is\_active : bool}
        \attr{created\_at : datetime}
    \end{class}

    \begin{class}[text width=3.7cm]{Enrollment}{5.6,0}
        \attr{id : int}
        \attr{course\_id : int}
        \attr{student\_id : int}
        \attr{status : EnrollmentStatus}
        \attr{requested\_at : datetime}
        \attr{decided\_at : datetime?}
    \end{class}

    \begin{class}[text width=3.7cm]{Course}{10.95,0}
        \attr{id : int}
        \attr{owner\_id : int}
        \attr{title : str}
        \attr{description : str}
        \attr{introduction : JSON?}
        \attr{summary : JSON?}
        \attr{is\_published : bool}
        \attr{is\_archived : bool}
        \attr{created\_at : datetime}
        \attr{updated\_at : datetime}
    \end{class}

    \begin{class}[text width=3.7cm]{LessonCompletion}{5.6,-5.7}
        \attr{user\_id : int}
        \attr{lesson\_id : int}
        \attr{passed\_at : datetime}
    \end{class}

    \begin{class}[text width=3.7cm]{Lesson}{10.95,-5.7}
        \attr{id : int}
        \attr{course\_id : int}
        \attr{title : str}
        \attr{position : int}
        \attr{outline : JSON?}
        \attr{conclusions : JSON?}
        \attr{is\_published : bool}
        \attr{is\_archived : bool}
        \attr{created\_at : datetime}
        \attr{updated\_at : datetime}
    \end{class}

    % User - Enrollment - Course (student)
    \coordinate (r1) at (0,-1.3);
    \draw (User.east |- r1) -- node[rolelabel, below] {uczeń} (Enrollment.west |- r1);
    \node[mult, above right] at (User.east |- r1) {1};
    \node[mult, above left] at (Enrollment.west |- r1) {*};
    \draw (Enrollment.east |- r1) -- (Course.west |- r1);
    \node[mult, above right] at (Enrollment.east |- r1) {*};
    \node[mult, above left] at (Course.west |- r1) {1};

    % User - Course (owner)
    \draw (User.north) -- ++(0,0.45) -| node[rolelabel, above, pos=0.25] {autor (owner\_id)} (Course.north);
    \node[mult, above right] at (User.north) {1};
    \node[mult, above right] at (Course.north) {*};

    % Course - Lesson (composition)
    \draw[comp] (Course.south) -- (Lesson.north);
    \node[mult, right=1.5mm] at ($(Course.south)+(0,-0.25)$) {1};
    \node[mult, above right] at (Lesson.north) {*};

    % LessonCompletion - Lesson
    \coordinate (r2) at (0,-6.6);
    \draw (LessonCompletion.east |- r2) -- (Lesson.west |- r2);
    \node[mult, above right] at (LessonCompletion.east |- r2) {*};
    \node[mult, above left] at (Lesson.west |- r2) {1};

    % User - LessonCompletion
    \draw (User.south) |- node[rolelabel, above, pos=0.75] {zaliczenia} (LessonCompletion.west);
    \node[mult, below right] at (User.south) {1};
    \node[mult, below left] at (LessonCompletion.west) {*};
\end{tikzpicture}

    \caption{Diagram klas: konta, kursy, zapisy i zaliczenia (\texttt{identity}, \texttt{curriculum})}
    \label{fig:classes-courses}
\end{figure}

\begin{figure}[H]
    \centering
    % Class diagram 2: lesson versions, files, quiz pool and drawn sets (content, assessment)
\begin{tikzpicture}[font=\scriptsize]
    \begin{class}[text width=2.4cm]{Lesson}{5.4,0}
    \end{class}

    \begin{class}[text width=3.7cm]{LessonContent}{0,-2.0}
        \attr{id : int}
        \attr{lesson\_id : int}
        \attr{form : V | A | K}
        \attr{blocks : JSON}
        \attr{updated\_at : datetime}
    \end{class}

    \begin{class}[text width=3.7cm]{LessonQuiz}{5.4,-2.0}
        \attr{id : int}
        \attr{lesson\_id : int}
        \attr{questions : JSON}
        \attr{updated\_at : datetime}
    \end{class}

    \begin{class}[text width=3.7cm]{QuizDraw}{10.8,-2.0}
        \attr{id : int}
        \attr{user\_id : int}
        \attr{lesson\_id : int}
        \attr{question\_ids : JSON}
        \attr{answers : JSON?}
        \attr{percent : int?}
        \attr{form : V | A | K ?}
        \attr{created\_at : datetime}
        \attr{submitted\_at : datetime?}
    \end{class}

    \begin{class}[text width=3.7cm]{File}{0,-5.1}
        \attr{id : int}
        \attr{owner\_id : int}
        \attr{original\_name : str}
        \attr{content\_type : str}
        \attr{size : int}
        \attr{storage\_name : str}
        \attr{created\_at : datetime}
    \end{class}

    \begin{class}[text width=2.4cm]{User}{10.8,-6.65}
    \end{class}

    % Lesson - LessonContent (three versions)
    \draw[comp] (Lesson.west) -| (LessonContent.north);
    \node[mult, below left] at ($(Lesson.west)+(-0.3,0)$) {1};
    \node[mult, above left] at (LessonContent.north) {0..3};
    \node[rolelabel, above] at ($(Lesson.west)!0.5!(LessonContent.north |- Lesson.west)$) {wersje V, A, K};

    % Lesson - LessonQuiz (pool)
    \draw[comp] (Lesson.south) -- (LessonQuiz.north);
    \node[mult, right=1.5mm] at ($(Lesson.south)+(0,-0.3)$) {1};
    \node[mult, above right] at (LessonQuiz.north) {0..1};
    \node[rolelabel, left=1.5mm] at ($(Lesson.south)!0.5!(LessonQuiz.north)$) {pula};

    % Lesson - QuizDraw
    \draw (Lesson.east) -| (QuizDraw.north);
    \node[mult, above right] at (Lesson.east) {1};
    \node[mult, above right] at (QuizDraw.north) {*};

    % QuizDraw ..> LessonQuiz (ids of drawn questions)
    \coordinate (r1) at (0,-2.9);
    \draw[dep] (QuizDraw.west |- r1) -- (LessonQuiz.east |- r1);
    \node[rolelabel, below] at ($(QuizDraw.west |- r1)!0.5!(LessonQuiz.east |- r1)$) {5 z puli};

    % LessonContent ..> File
    \draw[dep] (LessonContent.south) -- node[rolelabel, right] {bloki: file\_id} (File.north);

    % User - QuizDraw
    \draw (User.north) -- (QuizDraw.south);
    \node[mult, above right] at (User.north) {1};
    \node[mult, below right] at (QuizDraw.south) {*};

    % User - File (owner)
    \draw (File.east |- User.west) -- node[rolelabel, above] {autor (owner\_id)} (User.west);
    \node[mult, above right] at (File.east |- User.west) {*};
    \node[mult, above left] at (User.west) {1};
\end{tikzpicture}

    \caption{Diagram klas: wersje lekcji, pliki, pula i zestawy pytań (\texttt{content}, \texttt{assessment})}
    \label{fig:classes-content}
\end{figure}

\begin{figure}[H]
    \centering
    % Class diagram 3: student model and decision log (learning_style)
\begin{tikzpicture}[font=\scriptsize]
    % Top row
    \begin{class}[text width=3.45cm]{SurveyResult}{0,0}
        \attr{id : int}
        \attr{user\_id : int}
        \attr{points\_v : int}
        \attr{points\_a : int}
        \attr{points\_k : int}
        \attr{dominant\_style : str}
        \attr{created\_at : datetime}
    \end{class}

    \begin{class}[text width=3.45cm]{LearningProfile}{4.05,0}
        \attr{user\_id : int \{PK\}}
        \attr{style : V | A | K ?}
        \attr{style\_source : str}
        \attr{percent\_v : int}
        \attr{percent\_a : int}
        \attr{percent\_k : int}
        \attr{updated\_at : datetime}
    \end{class}

    \begin{class}[text width=3.45cm]{StyleHistory}{8.1,0}
        \attr{id : int}
        \attr{user\_id : int}
        \attr{percent\_v : int}
        \attr{percent\_a : int}
        \attr{percent\_k : int}
        \attr{style : V | A | K ?}
        \attr{event\_id : int?}
        \attr{created\_at : datetime}
    \end{class}

    \begin{class}[text width=3.45cm]{DecisionEvent}{12.15,0}
        \attr{id : int}
        \attr{user\_id : int}
        \attr{type : str}
        \attr{course\_id : int?}
        \attr{lesson\_id : int?}
        \attr{data : JSON}
        \attr{created\_at : datetime}
    \end{class}

    % The student: every class above and below refers to it
    \node[draw, fill=white, minimum height=0.6cm, text width=15.5cm, align=center,
          font=\scriptsize\bfseries] (User) at (6.075,-4.25) {User};

    % Bottom row
    \begin{class}[text width=3.45cm]{QuizResult}{0,-5.25}
        \attr{id : int}
        \attr{user\_id : int}
        \attr{course\_id : int}
        \attr{lesson\_id : int}
        \attr{draw\_id : int?}
        \attr{form : V | A | K}
        \attr{percent : int}
        \attr{passed : bool}
        \attr{counted : bool}
        \attr{created\_at : datetime}
    \end{class}

    \begin{class}[text width=3.45cm]{LessonForm}{4.05,-5.25}
        \attr{user\_id : int \{PK\}}
        \attr{lesson\_id : int \{PK\}}
        \attr{form : V | A | K}
        \attr{updated\_at : datetime}
    \end{class}

    \begin{class}[text width=3.6cm]{StyleProposal}{8.1,-5.25}
        \attr{id : int}
        \attr{user\_id : int}
        \attr{kind : str}
        \attr{from\_style : str}
        \attr{to\_style : str}
        \attr{evidence : JSON}
        \attr{status : str}
        \attr{counted\_at\_decision : int?}
        \attr{created\_at : datetime}
        \attr{decided\_at : datetime?}
    \end{class}

    \begin{class}[text width=3.45cm]{StyleTrial}{12.15,-5.25}
        \attr{id : int}
        \attr{user\_id : int}
        \attr{from\_style : str}
        \attr{to\_style : str}
        \attr{source : str}
        \attr{before\_average : int?}
        \attr{after\_average : int?}
        \attr{result : str?}
        \attr{started\_at : datetime}
        \attr{finished\_at : datetime?}
    \end{class}

    % Survey catalogue
    \begin{class}[text width=3.45cm]{SurveyQuestion}{0,-10.1}
        \attr{id : int}
        \attr{text : str}
        \attr{position : int}
        \attr{is\_active : bool}
    \end{class}

    \begin{class}[text width=3.45cm]{SurveyAnswer}{5.4,-10.1}
        \attr{id : int}
        \attr{question\_id : int}
        \attr{text : str}
        \attr{style : V | A | K}
        \attr{position : int}
    \end{class}

    % Associations with User
    \foreach \cls/\m in {SurveyResult/*, LearningProfile/0..1, StyleHistory/*, DecisionEvent/*} {
        \draw (\cls.south) -- (\cls.south |- User.north);
        \node[mult, below right] at (\cls.south) {\m};
        \node[mult, above right] at (\cls.south |- User.north) {1};
    }
    \foreach \cls in {QuizResult, LessonForm, StyleProposal, StyleTrial} {
        \draw (\cls.north) -- (\cls.north |- User.south);
        \node[mult, above right] at (\cls.north) {*};
        \node[mult, below right] at (\cls.north |- User.south) {1};
    }

    % StyleHistory ..> DecisionEvent (event_id)
    \coordinate (r1) at (0,-1.0);
    \draw[dep] (StyleHistory.east |- r1) -- (DecisionEvent.west |- r1);

    % SurveyQuestion - SurveyAnswer
    \coordinate (r2) at (0,-11.0);
    \draw[comp] (SurveyQuestion.east |- r2) -- (SurveyAnswer.west |- r2);
    \node[mult, below right] at ($(SurveyQuestion.east |- r2)+(0.25,0)$) {1};
    \node[mult, below left] at (SurveyAnswer.west |- r2) {3};
\end{tikzpicture}

    \caption{Diagram klas: model ucznia i dziennik decyzji (\texttt{learning\_style})}
    \label{fig:classes-student}
\end{figure}

\subsection{Uporządkowany alfabetycznie wykaz klas}

\textbf{Course} -- kurs nauczyciela: opis w katalogu, wstęp, ,,Podsumowanie i wnioski'' oraz lekcje. Kurs jest prywatny albo opublikowany (opublikować można kurs z co najmniej jedną opublikowaną lekcją); usunięcie oznacza archiwizację. \\
\textit{Atrybuty:} \texttt{id}, \texttt{owner\_id}, \texttt{title}, \texttt{description}, \texttt{introduction}, \texttt{summary} (dokumenty edytora), \texttt{is\_published}, \texttt{is\_archived}, \texttt{created\_at}, \texttt{updated\_at}.

\vspace{0.5em}
\textbf{DecisionEvent} -- wpis dziennika decyzji: typ zdarzenia i dane (liczby), z których decyzja wynika; kurs i lekcja, gdy zdarzenie ich dotyczy. \\
\textit{Atrybuty:} \texttt{id}, \texttt{user\_id}, \texttt{type}, \texttt{course\_id}, \texttt{lesson\_id}, \texttt{data} (JSON), \texttt{created\_at}.

\vspace{0.5em}
\textbf{Enrollment} -- prośba ucznia o zapis na kurs i decyzja autora: \texttt{pending}, \texttt{accepted} albo \texttt{rejected}; jedna para (kurs, uczeń). \\
\textit{Atrybuty:} \texttt{id}, \texttt{course\_id}, \texttt{student\_id}, \texttt{status}, \texttt{requested\_at}, \texttt{decided\_at}.

\vspace{0.5em}
\textbf{File} -- plik wgrany przez nauczyciela (obraz, nagranie, film, PDF), przechowywany na dysku pod losową nazwą; bloki lekcji wskazują go przez \texttt{file\_id}. \\
\textit{Atrybuty:} \texttt{id}, \texttt{owner\_id}, \texttt{original\_name}, \texttt{content\_type}, \texttt{size}, \texttt{storage\_name}, \texttt{created\_at}.

\vspace{0.5em}
\textbf{LearningProfile} -- profil ucznia: jeden aktywny styl (\texttt{V}, \texttt{A}, \texttt{K}; \texttt{null} = nierozstrzygnięty remis), źródło stylu (\texttt{survey}, \texttt{tie\_choice}, \texttt{manual}, \texttt{proposal}) i procenty form. Ograniczenie w bazie: każdy procent 5--90, suma 100. Brak wiersza oznacza, że uczeń nie wypełnił ankiety. \\
\textit{Atrybuty:} \texttt{user\_id}, \texttt{style}, \texttt{style\_source}, \texttt{percent\_v}, \texttt{percent\_a}, \texttt{percent\_k}, \texttt{updated\_at}.

\vspace{0.5em}
\textbf{Lesson} -- lekcja kursu w kolejności nauczyciela: zarys, wnioski i stan publikacji. \\
\textit{Atrybuty:} \texttt{id}, \texttt{course\_id}, \texttt{title}, \texttt{position}, \texttt{outline}, \texttt{conclusions}, \texttt{is\_published}, \texttt{is\_archived}, \texttt{created\_at}, \texttt{updated\_at}.

\vspace{0.5em}
\textbf{LessonCompletion} -- pierwsze zaliczenie lekcji przez ucznia; podstawa odblokowania następnej lekcji (R10). \\
\textit{Atrybuty:} \texttt{user\_id}, \texttt{lesson\_id}, \texttt{passed\_at}.

\vspace{0.5em}
\textbf{LessonContent} -- jedna wersja lekcji (V, A albo K) jako uporządkowana lista bloków: tekst oraz załączniki (obraz, film, nagranie, PDF, link). Para (\texttt{lesson\_id}, \texttt{form}) jest unikalna. \\
\textit{Atrybuty:} \texttt{id}, \texttt{lesson\_id}, \texttt{form}, \texttt{blocks} (JSON), \texttt{updated\_at}.

\vspace{0.5em}
\textbf{LessonForm} -- jawny wybór formy lekcji przez ucznia (,,Przerób tę lekcję w innej formie''); brak wiersza oznacza formę równą aktywnemu stylowi. \\
\textit{Atrybuty:} \texttt{user\_id}, \texttt{lesson\_id}, \texttt{form}, \texttt{updated\_at}.

\vspace{0.5em}
\textbf{LessonQuiz} -- pula pytań lekcji: 1--50 pytań jednokrotnego wyboru (3--4 odpowiedzi, jedna poprawna, wyjaśnienie 1--1000 znaków); do publikacji lekcji potrzeba co najmniej 5 pytań. \\
\textit{Atrybuty:} \texttt{id}, \texttt{lesson\_id}, \texttt{questions} (JSON), \texttt{updated\_at}.

\vspace{0.5em}
\textbf{QuizDraw} -- zestaw pytań wylosowany dla jednego podejścia (identyfikatory w kolejności pokazania); ocena tylko wobec tego zestawu. Zestaw bez \texttt{submitted\_at} jest otwarty i wraca po odświeżeniu strony. \\
\textit{Atrybuty:} \texttt{id}, \texttt{user\_id}, \texttt{lesson\_id}, \texttt{question\_ids}, \texttt{answers}, \texttt{percent}, \texttt{created\_at}, \texttt{submitted\_at}.

\vspace{0.5em}
\textbf{QuizResult} -- wynik podejścia w modelu ucznia: forma lekcji w chwili wysłania, procent, zaliczenie i to, czy wynik zmienił procenty (\texttt{counted}). \\
\textit{Atrybuty:} \texttt{id}, \texttt{user\_id}, \texttt{course\_id}, \texttt{lesson\_id}, \texttt{draw\_id}, \texttt{form}, \texttt{percent}, \texttt{passed}, \texttt{counted}, \texttt{created\_at}.

\vspace{0.5em}
\textbf{StyleHistory} -- stan profilu po każdej zmianie (ankieta, policzony wynik, zmiana stylu) z odwołaniem do zdarzenia dziennika; dane wykresu ,,Procenty form w czasie''. \\
\textit{Atrybuty:} \texttt{id}, \texttt{user\_id}, \texttt{percent\_v}, \texttt{percent\_a}, \texttt{percent\_k}, \texttt{style}, \texttt{event\_id}, \texttt{created\_at}.

\vspace{0.5em}
\textbf{StyleProposal} -- propozycja zmiany stylu: \texttt{switch} (na formę potwierdzoną wynikami) albo \texttt{revert} (powrót po nieudanym okresie próbnym), z dowodem w \texttt{evidence}. Stan: \texttt{open}, \texttt{accepted}, \texttt{rejected}, \texttt{expired}, \texttt{superseded}; najwyżej jedna otwarta na ucznia (częściowy indeks unikalny). \\
\textit{Atrybuty:} \texttt{id}, \texttt{user\_id}, \texttt{kind}, \texttt{from\_style}, \texttt{to\_style}, \texttt{evidence}, \texttt{status}, \texttt{counted\_at\_decision}, \texttt{created\_at}, \texttt{decided\_at}.

\vspace{0.5em}
\textbf{StyleTrial} -- okres próbny po zmianie stylu (\texttt{proposal} albo \texttt{manual}): średnia przed i po oraz wynik \texttt{helped}, \texttt{neutral}, \texttt{worse} albo \texttt{no\_baseline}. \\
\textit{Atrybuty:} \texttt{id}, \texttt{user\_id}, \texttt{from\_style}, \texttt{to\_style}, \texttt{source}, \texttt{before\_average}, \texttt{after\_average}, \texttt{result}, \texttt{started\_at}, \texttt{finished\_at}.

\vspace{0.5em}
\textbf{SurveyAnswer} -- odpowiedź na pytanie ankiety, przypisana do jednej formy V, A lub K. \\
\textit{Atrybuty:} \texttt{id}, \texttt{question\_id}, \texttt{text}, \texttt{style}, \texttt{position}.

\vspace{0.5em}
\textbf{SurveyQuestion} -- pytanie ankiety VAK (20 pytań, po 3 odpowiedzi). \\
\textit{Atrybuty:} \texttt{id}, \texttt{text}, \texttt{position}, \texttt{is\_active}.

\vspace{0.5em}
\textbf{SurveyResult} -- wynik jedynego wypełnienia ankiety; przy remisie \texttt{dominant\_style} zapisuje litery remisu (np. \texttt{VA}). \\
\textit{Atrybuty:} \texttt{id}, \texttt{user\_id}, \texttt{points\_v}, \texttt{points\_a}, \texttt{points\_k}, \texttt{dominant\_style}, \texttt{created\_at}.

\vspace{0.5em}
\textbf{User} -- konto z rolą \texttt{student}, \texttt{teacher} albo \texttt{admin}. Stan zatwierdzenia (\texttt{approved}, \texttt{pending}, \texttt{rejected}) dotyczy nauczycieli z rejestracji; konto można zablokować. \\
\textit{Atrybuty:} \texttt{id}, \texttt{email}, \texttt{password\_hash}, \texttt{display\_name}, \texttt{role}, \texttt{approval\_status}, \texttt{is\_active}, \texttt{created\_at}.

\subsection{Model danych}

Baza danych PostgreSQL zawiera 19 tabel odpowiadających klasom z diagramu (schemat tworzą migracje Alembic; v3 wprowadza migracja \texttt{01324f58656d\_v3\_adaptation}). Najważniejsze relacje:

\begin{itemize}[leftmargin=*]
    \item \texttt{users} 1:N \texttt{courses} (autor), \texttt{enrollments}, \texttt{lesson\_completions}, \texttt{files}, \texttt{quiz\_draws}; 1:1 \texttt{learning\_profiles}; 1:N \texttt{survey\_results}, \texttt{style\_history}, \texttt{quiz\_results}, \texttt{lesson\_forms}, \texttt{style\_proposals}, \texttt{style\_trials}, \texttt{decision\_events}
    \item \texttt{courses} 1:N \texttt{lessons}; \texttt{courses} N:M \texttt{users} przez \texttt{enrollments} (unikalna para kurs--uczeń)
    \item \texttt{lessons} 1:0..3 \texttt{lesson\_contents} (unikalne \texttt{lesson\_id}, \texttt{form}), 1:0..1 \texttt{lesson\_quizzes}, 1:N \texttt{quiz\_draws}
    \item \texttt{quiz\_results.draw\_id} $\to$ \texttt{quiz\_draws}, \texttt{style\_history.event\_id} $\to$ \texttt{decision\_events} (\texttt{SET NULL}); dane ucznia są usuwane kaskadowo razem z kontem
\end{itemize}

Klucze obce między modułami istnieją tylko w bazie -- w kodzie moduły korzystają ze siebie przez funkcje serwisów (rozdział~\ref{sec:architecture}).

\subsection{Reguły adaptacji}
\label{sec:rules}

Wszystkie reguły adaptacji są czystymi funkcjami w jednym module \nolinkurl{backend/app/modules/learning_style/rules.py} (bez SQLAlchemy i FastAPI): wejściem są proste struktury, wyjściem decyzja z danymi uzasadnienia. Losowanie pytań jest w \nolinkurl{assessment/drawing.py}, blokada lekcji w \nolinkurl{curriculum/progress.py}, a braki do publikacji w plikach \texttt{readiness.py} modułów \texttt{curriculum}, \texttt{content} i \texttt{assessment}. Reguły nie wymagają danych treningowych ani zewnętrznego modelu AI; każda reakcja niesie liczby, z których wynika.

\begin{table}[H]
\centering
\small
\begin{tabular}{|l|c|p{8.2cm}|}
\hline
\textbf{Stała} & \textbf{Wartość} & \textbf{Znaczenie} \\ \hline
\texttt{PASS\_THRESHOLD} & 70 & zaliczenie: wynik $\geq$ 70\% (przy 5 pytaniach: od 4 z 5) \\ \hline
\texttt{STEP} & 10 & zmiana procentu formy po policzonym wyniku \\ \hline
\texttt{MIN\_PERCENT} & 5 & minimum każdej formy, więc maksimum 90 \\ \hline
\texttt{QUIZ\_SIZE} & 5 & pytań na podejście; pula do publikacji $\geq$ 5 \\ \hline
\texttt{TRIAL\_LESSONS} & 2 & długość okresu próbnego (policzone wyniki) \\ \hline
\texttt{HELPED\_GAIN} & 10 & wzrost średniej, od którego zmiana ,,pomogła'' \\ \hline
\texttt{REVERT\_DROP} & 10 & spadek średniej, od którego proponujemy powrót \\ \hline
\texttt{RENEW\_AFTER} & 2 & nowe policzone wyniki, po których wraca odrzucona propozycja \\ \hline
\texttt{HINT\_MIN\_STUDENTS} & 2 & minimum uczniów w formie dla podpowiedzi z kursu \\ \hline
\end{tabular}
\caption{Stałe reguł adaptacji (\texttt{rules.py}); remisy rozstrzyga kolejność V, A, K}
\end{table}

\subsubsection{R1. Profil z ankiety}
\texttt{survey\_percentages(points)}: udział formy to $p_f / \sum p \cdot 100$; forma poniżej 5\% dostaje 5\%, a pozostałe dzielą resztę proporcjonalnie do punktów. Zaokrąglenie \textbf{metodą największych reszt} (podłoga, brakujące jednostki dla największych części ułamkowych), więc suma zawsze wynosi 100. Przykłady: A 11 / V 6 / K 3 $\to$ \textbf{A 55 / V 30 / K 15}; V 9 / A 8 / K 3 $\to$ V 45 / A 40 / K 15; V 20 / A 0 / K 0 $\to$ V 90 / A 5 / K 5.

\texttt{survey\_leaders(points)}: formy z maksimum punktów. Jedna forma -- to jest styl (\texttt{style\_source = survey}). Kilka -- profil bez stylu, uczeń wybiera jeden z remisujących (\texttt{tie\_choice}); do wyboru kursy są niedostępne. Ankietę wypełnia się raz.

\subsubsection{R2. Zmiana procentów po quizie}
\texttt{apply\_result(percentages, form, passed)}: procent formy lekcji zmienia się o $\pm 10$ (w granicach 5--90), a dwie pozostałe formy dzielą $100 - \textit{nowy}$ proporcjonalnie do swoich poprzednich wartości (forma poniżej 5 dostaje 5); zaokrąglenie jak w R1. Przykłady sprawdzane w testach:
\begin{itemize}[leftmargin=*]
    \item V 45 / A 40 / K 15, niezaliczone w V $\to$ \textbf{V 35 / A 47 / K 18} (A $47{,}27 \to 47$, K $17{,}73 \to 18$)
    \item V 35 / A 47 / K 18, zaliczone w A $\to$ \textbf{V 28 / A 57 / K 15} (V $28{,}40 \to 28$, K $14{,}60 \to 15$)
    \item V 85 / A 10 / K 5, zaliczone w V $\to$ V 90 / A 5 / K 5 (limit 90 i minimum 5)
\end{itemize}

\texttt{is\_counted(previous\_forms\_in\_lesson, form)}: procenty zmienia \textbf{tylko pierwsze podejście w danej formie w danej lekcji} (zaliczone albo nie). Kolejne podejścia w tej samej formie są zapisywane (\texttt{counted = false}) i mogą zaliczyć lekcję, ale procentów nie zmieniają -- nie da się ich ,,nabić''.

\subsubsection{R3. Statystyki form ucznia}
\texttt{form\_stats(results)}: dla każdej formy liczba policzonych wyników (\texttt{lessons}) i ich średnia (\texttt{average}, \texttt{null} bez danych) ze wszystkich kursów ucznia. Forma jest \textbf{potwierdzona wynikami} (\texttt{confirmed}), gdy uczeń ma w niej co najmniej jedno zaliczone podejście (także niepoliczone); inaczej jest ,,tylko z ankiety''.

\subsubsection{R4. Zaliczenie: odblokowanie i pochwała}
Każde zaliczone podejście zapisuje \texttt{lesson\_completions} (raz) i odblokowuje następną opublikowaną lekcję. \texttt{praise(results)}: gdy bieżący wynik jest zaliczony i policzony, a końcowa seria policzonych zaliczeń w tej samej formie ma co najmniej 2 lekcje, uczeń dostaje pochwałę z danymi, np. ,,3 lekcje z rzędu zaliczone w formie wzrokowej, średnio 87\%'' (80, 100, 80).

\subsubsection{R5. Niezaliczenie: sugestia ,,przerób w formie X''}
\texttt{retry\_suggestion(...)}: kandydaci to dwie formy różne od formy, w której uczeń przerabiał lekcję.
\begin{itemize}[leftmargin=*]
    \item Jeśli któryś kandydat jest potwierdzony (R3), X to potwierdzony kandydat z najwyższym procentem (remis $\to$ wyższa średnia $\to$ V, A, K); powód \texttt{confirmed}, dowód: najlepsze zaliczenie w X (lekcja, wynik)
    \item W przeciwnym razie X to kandydat z najwyższym procentem \textbf{w ankiecie} (remis $\to$ wyższy bieżący procent $\to$ V, A, K); powód \texttt{survey}, dowód: procent X w ankiecie
\end{itemize}
Do sugestii dołączony jest wynik oblanej formy w tym quizie i poprzedni policzony wynik w tej formie w innej lekcji. Przykład (Kuba, lekcja 3): ,,Wzrokowe: ten quiz 40\%, poprzednia lekcja 40\%. Słuchowe były drugie w Twojej ankiecie (40\%), nie mamy jeszcze o nich wyników -- warto sprawdzić''. Sugestia trafia do dziennika przy każdym niezaliczonym podejściu; zmiana formy dotyczy tylko tej lekcji.

\subsubsection{R6. Propozycja zmiany stylu na stałe}
\texttt{switch\_proposal(style, percentages, stats, rejections, trial\_active)}: system proponuje formę Y różną od stylu, gdy Y jest \textbf{potwierdzona} wynikami i ma \textbf{wyższy procent} niż aktywny styl (kilka kandydatów: najwyższy procent $\to$ wyższa średnia $\to$ V, A, K). Forma ,,tylko z ankiety'' nigdy nie dostaje propozycji. Propozycji nie ma w trakcie okresu próbnego ani dla formy odrzuconej, dopóki od odrzucenia nie przybędą 2 nowe policzone wyniki. Otwarta propozycja, której warunek przestał obowiązywać, wygasa (\texttt{expired}); ręczna zmiana stylu zamyka ją jako \texttt{superseded}. Dowód zawiera liczbę lekcji, średnią i procent obu form, np. ,,Wzrokowe: 3 lekcje, średnio 53\% · Słuchowe: 1 lekcja, średnio 80\%''.

\texttt{forms\_to\_clear(chosen\_forms, passed\_lesson\_ids)}: przy \textbf{każdej} zmianie stylu (przyjęta propozycja \texttt{switch} lub \texttt{revert}, zmiana ręczna) usuwane są jawne wybory formy w lekcjach jeszcze niezaliczonych -- nowy styl obowiązuje wszędzie poza zaliczonymi lekcjami. Lista wyczyszczonych lekcji trafia do danych zdarzenia.

\subsubsection{R7. Okres próbny}
Zaczyna się przy przyjęciu propozycji \texttt{switch} albo ręcznej zmianie stylu (po przyjęciu propozycji powrotu -- nie, żeby nie było pętli). \texttt{trial\_baseline}: średnia ,,przed'' to średnia z ostatnich co najwyżej 2 policzonych wyników w poprzednim stylu (\texttt{null}, gdy brak). \texttt{trial\_outcome}: po 2 policzonych wynikach w nowym stylu liczona jest średnia ,,po'' i różnica $d = \textit{po} - \textit{przed}$:

\begin{table}[H]
\centering
\begin{tabular}{|l|l|p{7.2cm}|}
\hline
\textbf{Warunek} & \textbf{Wynik} & \textbf{Reakcja systemu} \\ \hline
$d \geq +10$ & \texttt{helped} & ,,Zmiana pomogła: średnia 40\% $\to$ 90\%'' \\ \hline
$-10 < d < +10$ & \texttt{neutral} & ,,Bez wyraźnej różnicy: średnia 60\% $\to$ 65\%'' \\ \hline
$d \leq -10$ & \texttt{worse} & propozycja powrotu (\texttt{revert}) do poprzedniego stylu, bez wymogu potwierdzenia \\ \hline
brak ,,przed'' & \texttt{no\_baseline} & ,,Średnia w nowej formie: 85\%'' \\ \hline
\end{tabular}
\caption{Wynik okresu próbnego (granice sprawdzane w testach: $+10$, $+9$, $-9$, $-10$)}
\end{table}

\subsubsection{R8. Podpowiedź z kursu}
\texttt{course\_hint(course\_results, student\_id)}: tylko dla ucznia \textbf{bez wyniku w kursie}. Policzone wyniki innych uczniów kursu grupowane są po formie; brane są formy z co najmniej 2 uczniami; wygrywa najwyższa średnia (remis $\to$ więcej uczniów $\to$ V, A, K), np. ,,W tym kursie najlepiej wypada forma słuchowa (średnio 86\%, 3 uczniów)''. To tylko informacja -- system niczego nie zmienia.

\subsubsection{R9. Losowanie pytań}
\texttt{draw\_questions(pool\_ids, last\_ids, k=5, rng)}: najpierw losowo pytania spoza ostatniego wysłanego zestawu, potem uzupełnienie z ostatniego, razem $\min(5, |\textit{pula}|)$. Przy puli 8 powtórka daje 3 nowe pytania i 2 z poprzedniego zestawu. Otwarty (niewysłany) zestaw jest zwracany ponownie, więc odświeżenie strony nie losuje od nowa.

\subsubsection{R10. Blokada lekcji}
\texttt{lesson\_states(published\_lessons, passed\_ids)}: lekcja jest \texttt{passed}, gdy uczeń ją zaliczył; \texttt{available}, gdy jest pierwsza albo poprzednia opublikowana jest zaliczona; inaczej \texttt{locked} z numerem wymaganej lekcji. Zaliczona lekcja zostaje dostępna także po zmianie kolejności. Serwer sprawdza to przy każdym dostępie ucznia do treści i quizu (\texttt{403 LESSON\_LOCKED}).

\subsubsection{R11. Gotowość lekcji do publikacji}
Braki zgłaszane kodami: \texttt{OUTLINE\_MISSING}, \texttt{CONCLUSIONS\_MISSING} (\texttt{curriculum}); dla każdej wersji \texttt{VERSION\_EMPTY}, a ponadto V -- \texttt{V\_VISUAL\_MISSING} (tekst i obraz lub film), A -- \texttt{A\_AUDIO\_MISSING} i \texttt{A\_DESCRIPTION\_MISSING} (nagranie i opis), K -- \texttt{K\_INSTRUCTION\_MISSING} i \texttt{K\_WORKSHEET\_MISSING} (instrukcja i karta PDF) (\texttt{content}); \texttt{QUIZ\_TOO\_FEW\_QUESTIONS} z liczbą pytań i minimum 5 (\texttt{assessment}). Opublikowana lekcja nie może stać się niekompletna (\texttt{409 PUBLISHED\_LESSON\_INCOMPLETE}).

\subsubsection{Przykład: ścieżka ucznia z danych demonstracyjnych}
\label{sec:rules-example}

Kuba (\texttt{uczen.demo@imsu.pl}) uzyskał w ankiecie V 9 / A 8 / K 3, więc jego styl to V, a profil V 45 / A 40 / K 15 (R1). Seed buduje jego historię w kursie ,,Matematyka: funkcja liniowa'' prawdziwymi serwisami (wiersze 1--3); wiersze 4--6 to przebieg pokazu na żywo (scenariusz \nolinkurl{docs/SCENARIUSZ-POKAZU.md}); po pokazie \texttt{-{}-reset-students} przywraca stan z wiersza 3 (rozdział~\ref{sec:setup}).

\begin{table}[H]
\centering
\small
\begin{tabular}{|l|c|c|c|c|p{5.4cm}|}
\hline
\textbf{Zdarzenie} & \textbf{Forma} & \textbf{Wynik} & \textbf{Policz.} & \textbf{V / A / K} & \textbf{Reakcja systemu} \\ \hline
Ankieta & -- & -- & -- & 45 / 40 / 15 & styl V \\ \hline
1. Lekcja 1 & V & 80\% & tak & 55 / 33 / 12 & lekcja 2 odblokowana \\ \hline
2. Lekcja 2 & V & 40\% & tak & 45 / 40 / 15 & sugestia A (z ankiety) -- zignorowana \\ \hline
3. Lekcja 2 & V & 80\% & nie & 45 / 40 / 15 & lekcja 3 odblokowana, procenty bez zmian \\ \hline
4. Lekcja 3 & V & 40\% & tak & 35 / 47 / 18 & sugestia A (ankieta 40\%) \\ \hline
5. Lekcja 3 & A & 80\% & tak & 28 / 57 / 15 & lekcja 4, propozycja V $\to$ A \\ \hline
6. Przyjęcie & -- & -- & -- & 28 / 57 / 15 & styl A, okres próbny, ,,przed'' 40\% \\ \hline
\end{tabular}
\caption{Ścieżka Kuby: procenty, sugestie i propozycja}
\end{table}

W wierszu 5 propozycja powstaje, bo A ma zaliczenie (potwierdzona, R3) i procent 57 wyższy niż V 28 (R6). Dowód: V -- 3 policzone lekcje (80, 40, 40), średnio 53\%; A -- 1 lekcja, 80\%. Średnia ,,przed'' okresu próbnego to średnia dwóch ostatnich policzonych wyników w V: (40 + 40) / 2 = 40\%. Jeśli dwie następne lekcje w A dadzą np. 80\% i 100\% (średnia 90\%), $d = +50$ i okres kończy się wynikiem ,,Zmiana pomogła''. Taki przebieg ma w danych demo Marta (\texttt{uczen.zmiana@imsu.pl}): ankieta K 10 / V 7 / A 3 $\to$ K 50 / V 35 / A 15, lekcja 1 w K 40\%, sugestia V, lekcja 1 w V 80\%, przyjęta propozycja V, potem 100\% i 80\% w V -- ,,Zmiana pomogła: średnia 40\% $\to$ 90\%''.

\subsection{Kolejność kroków po quizie}

Moduł \texttt{assessment} ocenia odpowiedzi na serwerze i wywołuje jedną funkcję \texttt{learning\_style.hooks.record\_quiz\_result} -- jedyny punkt styku quizu z modelem ucznia (żaden endpoint nie przyjmuje wyniku od klienta). Funkcja wykonuje kroki w ustalonej kolejności:

\begin{enumerate}[leftmargin=*]
    \item Forma = forma lekcji w chwili wysłania (\texttt{lesson\_forms} albo aktywny styl).
    \item \texttt{is\_counted}; zapis \texttt{QuizResult}; jeśli policzony -- \texttt{apply\_result} i nowe procenty profilu (R2).
    \item Zaliczone: \texttt{curriculum.mark\_lesson\_passed} (odblokowanie następnej lekcji, R10), informacja o ukończeniu kursu, \texttt{praise} (R4). Zdarzenie \texttt{quiz\_attempt} z procentami przed i po; jeśli wynik policzony -- punkt w \texttt{style\_history}.
    \item Niezaliczone: \texttt{retry\_suggestion} i zdarzenie \texttt{retry\_suggestion} (R5).
    \item Trwający okres próbny: \texttt{trial\_outcome}; po 2 wynikach zdarzenie \texttt{trial\_finished}, a przy \texttt{worse} -- propozycja powrotu (R7).
    \item \texttt{switch\_proposal}: wygaszenie nieaktualnej propozycji, utworzenie nowej (R6). Otwarta propozycja powrotu czeka na decyzję ucznia i nie jest zastępowana.
    \item Jeden \texttt{commit} dla wyniku, profilu i wszystkich decyzji; zwrot \texttt{QuizFeedback} (forma, \texttt{counted}, procenty przed i po, odblokowana lekcja, pochwała, sugestia, propozycja, okres próbny).
\end{enumerate}

\subsection{Dziennik decyzji}

Każda decyzja systemu zapisuje wiersz \texttt{decision\_events} z typem i danymi. Serwer zwraca \textbf{kody i liczby}, a polskie zdania składa w jednym miejscu klient (\nolinkurl{features/learning-style/explain.ts}) -- tak samo dla ucznia (,,Moje postępy'') i nauczyciela (szczegóły ucznia w panelu kursu).

\begingroup
\small
\begin{longtable}{p{4.3cm} p{10.4cm}}
\toprule
\textbf{Typ} & \textbf{Dane (\texttt{data})} \\
\midrule
\endhead
\texttt{survey\_completed} & punkty, procenty, styl (\texttt{null} przy remisie), style z remisu \\
\texttt{tie\_resolved} & wybrany styl, style z remisu \\
\texttt{quiz\_attempt} & \texttt{draw\_id}, forma, wynik, zaliczenie, \texttt{counted}, procenty przed i po, odblokowana lekcja \\
\texttt{retry\_suggestion} & oblana forma (wynik, poprzedni wynik w tej formie), sugerowana forma (powód, procent, procent w ankiecie, najlepsze zaliczenie) \\
\texttt{lesson\_form\_changed} & forma przed i po, źródło (\texttt{manual} / \texttt{suggestion}) \\
\texttt{proposal\_created} & id, rodzaj, styl z / na, dowód \\
\texttt{proposal\_accepted} & id, rodzaj, styl z / na, wyczyszczone wybory formy \\
\texttt{proposal\_rejected}, \texttt{proposal\_expired} & id, rodzaj, styl z / na \\
\texttt{manual\_style\_change} & styl z / na, wyczyszczone wybory formy, zamknięta propozycja \\
\texttt{trial\_started} & id okresu, styl z / na, źródło, średnia przed \\
\texttt{trial\_finished} & id okresu, styl z / na, średnie przed i po, wynik \\
\bottomrule
\end{longtable}
\endgroup

Zdarzenia lekcji i quizu mają \texttt{course\_id} i \texttt{lesson\_id}; ankieta i zmiany stylu są globalne. Przykładowy wpis dziennika Kuby po kroku 4: \texttt{quiz\_attempt} z danymi \texttt{form = V}, \texttt{percent = 40}, \texttt{counted = true}, procenty przed V 45 / A 40 / K 15 i po V 35 / A 47 / K 18.

\section{Technologie}

\subsection{Architektura}
\label{sec:architecture}

System jest aplikacją internetową złożoną z trzech części: klienta działającego w przeglądarce, serwera udostępniającego REST API oraz bazy danych. Serwer to \textbf{modularny monolit} -- jedna aplikacja podzielona na moduły domenowe, które korzystają ze siebie wyłącznie przez publiczne funkcje serwisów:

\begin{table}[H]
\centering
\begin{tabular}{|l|p{10cm}|}
\hline
\textbf{Moduł} & \textbf{Odpowiedzialność} \\ \hline
\texttt{core} & konfiguracja, sesja bazy danych, format błędów (\texttt{AppError}: kod i parametry) \\ \hline
\texttt{identity} & konta, role, logowanie, sesja, administracja kont, zatwierdzanie nauczycieli \\ \hline
\texttt{curriculum} & kursy, lekcje (zarys, wnioski), zapisy na kurs, zaliczenia i blokada kolejności (R10), zmiana właściciela kursu przez administratora (\texttt{course\_owner\_changes}) \\ \hline
\texttt{content} & trzy wersje lekcji (bloki tekstu i załączniki, nagrania z transkrypcją), pliki, braki wersji (R11) \\ \hline
\texttt{learning\_style} & ankieta, profil procentowy, reguły R1--R8 i R7a (\texttt{rules.py}), wersja lekcji dla ucznia i stany lekcji (przyciski ,,Rozpocznij'' / ,,Kontynuuj'' / ,,Powtórz''), propozycje, okres próbny, dziennik decyzji, statystyki \\ \hline
\texttt{assessment} & pula pytań, losowanie zestawu (R9), ocenianie, wywołanie \texttt{record\_quiz\_result} \\ \hline
\texttt{authoring} & gotowość lekcji (zbiera braki z trzech modułów) i publikacja lekcji \\ \hline
\end{tabular}
\caption{Moduły serwera}
\end{table}

Zależności między modułami są jednokierunkowe:
\begin{center}
\small\ttfamily
core $\leftarrow$ identity $\leftarrow$ curriculum $\leftarrow$ content $\leftarrow$ learning\_style $\leftarrow$ assessment $\leftarrow$ authoring
\end{center}
Moduł może importować tylko moduły po swojej lewej stronie. Odblokowanie lekcji należy do \texttt{curriculum} (\texttt{lesson\_completions}), więc \texttt{learning\_style} po zaliczeniu wywołuje \texttt{curriculum.service.mark\_lesson\_passed} bez zależności ,,w górę''. \texttt{authoring} jest na końcu łańcucha, bo do decyzji o publikacji potrzebuje braków z \texttt{curriculum}, \texttt{content} i \texttt{assessment}.

Reguły adaptacji, losowania, blokady i gotowości są czystymi funkcjami, niezależnymi od frameworka i bazy danych. Wynik quizu trafia do modelu ucznia tylko przez \texttt{learning\_style.hooks.record\_quiz\_result}, a klient nie ma możliwości samodzielnego zgłoszenia wyniku.

\subsection{Technologie oprogramowania}

\textbf{Klient:}
\begin{itemize}[leftmargin=*]
    \item \textbf{Next.js 16} i \textbf{React 19} -- aplikacja internetowa z routingiem opartym na katalogach
    \item \textbf{TypeScript} -- typowany JavaScript
    \item \textbf{Tailwind CSS 4} i \textbf{shadcn/ui} (Base UI) -- style i komponenty interfejsu
    \item \textbf{TanStack Query} -- pobieranie i buforowanie danych z serwera
    \item \textbf{TipTap} -- edytor tekstu lekcji, zarysów i wniosków (formatowanie, tabele, kod w linii i bloki kodu czcionką o stałej szerokości, wklejanie z Worda)
    \item \textbf{KaTeX} -- wyświetlanie wzorów matematycznych
    \item \textbf{React Hook Form}, \textbf{react-dropzone}, \textbf{Lucide}, \textbf{Sonner} -- formularze, wgrywanie plików, ikony, komunikaty
    \item Wykres procentów form w czasie -- własny komponent SVG (bez biblioteki wykresów): punkty dla ankiety i policzonych quizów, zmiany stylu jako pionowe znaczniki, pod spodem tabela tych samych danych z kolumną ,,Zdarzenie''
\end{itemize}

\textbf{Serwer:}
\begin{itemize}[leftmargin=*]
    \item \textbf{Python 3.12} i \textbf{FastAPI} -- REST API
    \item \textbf{SQLAlchemy 2} i \textbf{psycopg 3} -- dostęp do bazy danych
    \item \textbf{Alembic} -- migracje schematu bazy
    \item \textbf{Pydantic} -- walidacja danych wejściowych i wyjściowych
    \item \textbf{PyJWT} i \textbf{bcrypt} -- sesja i haszowanie haseł
\end{itemize}

\textbf{Baza danych:}
\begin{itemize}[leftmargin=*]
    \item \textbf{PostgreSQL 17} uruchamiany w kontenerze \textbf{Docker} (Docker Compose); kolumny \textbf{JSONB} dla bloków lekcji, puli pytań, zestawów i danych dziennika
    \item Pliki wgrane przez nauczycieli przechowywane na dysku serwera, poza repozytorium
\end{itemize}

\textbf{Narzędzia deweloperskie:}
\begin{itemize}[leftmargin=*]
    \item \textbf{Git} i \textbf{GitHub} -- kontrola wersji
    \item \textbf{pytest} -- testy serwera; \textbf{Ruff} -- lint i formatowanie kodu Python
    \item \textbf{ESLint} i \textbf{Prettier} -- lint i formatowanie kodu klienta
    \item \textbf{Playwright} -- przejście scenariusza pokazu w przeglądarce (skrypt \nolinkurl{frontend/e2e/demo-tour.mjs}, uruchamiany ręcznie)
    \item \textbf{npm} i \textbf{pip} -- menedżery pakietów
    \item Generatory zasobów demo (jednorazowe, w osobnym środowisku, poza zależnościami projektu): \textbf{Pillow} (obrazy PNG i karty PDF), \textbf{edge-tts} i \textbf{imageio-ffmpeg} (nagrania MP3 wersji A; kompletność czytania sprawdzana znacznikami przeczytanych słów)
\end{itemize}

\subsection{Bezpieczeństwo}

\begin{itemize}[leftmargin=*]
    \item Sesja jako token JWT w ciasteczku \texttt{HttpOnly} (niedostępnym z JavaScriptu), \texttt{SameSite=Lax}
    \item Rola, aktywność i zatwierdzenie konta odczytywane z bazy przy każdym żądaniu -- blokada konta działa natychmiast, a niezatwierdzony nauczyciel dostaje \texttt{403 TEACHER\_NOT\_APPROVED}
    \item Hasła haszowane algorytmem bcrypt
    \item Uprawnienia sprawdzane na serwerze; zasób niedostępny dla użytkownika zwraca 404, bez zdradzania jego istnienia; treść lekcji i quiz tylko dla ucznia przyjętego na kurs i z odblokowaną lekcją; zmiana właściciela kursu tylko dla administratora i tylko na zatwierdzonego nauczyciela, z trwałym śladem w \texttt{course\_owner\_changes}
    \item Wgrywane pliki: biała lista rozszerzeń (bez SVG i HTML), limit 50 MB, losowe nazwy na dysku
    \item Ankieta i quiz wysyłane do klienta bez klucza odpowiedzi; zestaw pytań zapisany na serwerze, punkty i wynik liczy serwer
    \item CORS tylko dla adresu klienta, sekrety wyłącznie w plikach \texttt{.env} poza repozytorium
\end{itemize}

\subsection{Licencje}

Wszystkie wykorzystane technologie są dostępne na licencjach otwartych, pozwalających na swobodne użycie:

\begin{itemize}[leftmargin=*]
    \item Next.js, React, Tailwind CSS, shadcn/ui, TanStack Query, TipTap, KaTeX -- MIT License
    \item FastAPI, SQLAlchemy, Alembic, Pydantic, pytest -- MIT License
    \item Playwright -- Apache License 2.0
    \item PostgreSQL -- PostgreSQL License (podobna do MIT)
\end{itemize}

\section{Dokumentacja testów}

\subsection{Cel i zakres testowania}

Celem testów jest sprawdzenie reguł systemu (profil z ankiety, zmiana procentów, sugestie, propozycje, okres próbny, losowanie pytań, blokada lekcji, gotowość do publikacji) oraz uprawnień i poprawności wszystkich endpointów serwera. Testy obejmują wszystkie moduły serwera: konta i zatwierdzanie nauczycieli, kursy, zapisy i blokadę, wersje lekcji i pliki, model ucznia, quizy, publikację lekcji, a także dane demonstracyjne.

\subsection{Strategia testowania}

\begin{center}
\small
\begin{tabular}{llll}
\toprule
\textbf{Poziom} & Jednostkowy & Integracyjny & Systemowy \\
\midrule
\textbf{Typ} & white-box & black-box HTTP & black-box UI \\
\textbf{Narzędzie} & pytest & pytest + TestClient & Playwright (ręcznie) \\
\textbf{Zakres} & czyste funkcje reguł & endpointy REST z bazą & scenariusz pokazu \\
\bottomrule
\end{tabular}
\end{center}

\bigskip

\textbf{Testy jednostkowe} sprawdzają czyste funkcje bez bazy danych i bez frameworka: reguły R1--R8 (\nolinkurl{learning_style/rules.py}, każda reguła z liczbami ze specyfikacji, także własności ,,suma = 100 i każda forma $\geq$ 5'' dla siatki profili), losowanie R9 (\nolinkurl{assessment/drawing.py}, deterministyczne dzięki wstrzykiwanemu \texttt{rng}), ocenianie, blokadę R10 (\nolinkurl{curriculum/progress.py}), gotowość R11 (\texttt{readiness.py}), punktację ankiety, statystyki kursu i hasła.

\textbf{Testy integracyjne} wysyłają żądania HTTP do aplikacji przez \texttt{TestClient} i korzystają z prawdziwej bazy PostgreSQL (osobna baza \texttt{imsu\_test}). Schemat bazy testowej jest budowany tymi samymi migracjami Alembic co baza aplikacji, a każdy test działa w transakcji wycofywanej po jego zakończeniu, więc testy są od siebie niezależne. Wspólne dane (uczeń, nauczyciel, administrator, kurs z opublikowanymi lekcjami, zapis ucznia, ankieta) dostarczają fixture'y pytest.

\textbf{Test systemowy} to skrypt Playwright \nolinkurl{frontend/e2e/demo-tour.mjs}: przechodzi scenariusz pokazu na wszystkich kontach demonstracyjnych, zapisuje zrzut każdego ekranu (szerokość komputera i wybrane ekrany na telefonie 390 px), zgłasza poziome przewijanie strony, błędy JavaScript i odpowiedzi 5xx oraz sprawdza warunki ekranów: zmianę właściciela Programowania na Annę i z powrotem (wynik sprawdzany przez API), przyciski ,,Rozpocznij'' / ,,Powtórz'' i kłódki, co najmniej 2 nagrania w wersji A, każde z transkrypcją (otwieraną z klawiatury), blok kodu czcionką o stałej szerokości oraz ścieżkę A $\to$ V $\to$ A ucznia \texttt{uczen3} (brak nowego okresu próbnego, wpis o cofnięciu, przekreślony znacznik na wykresie). Ostatnie przejście: \textbf{50 zrzutów} w \nolinkurl{docs/zrzuty-v3}, bez błędów 500 i JS i bez poziomego przewijania.

\textbf{Test komponentu klienta} \nolinkurl{frontend/src/features/progress/chart-data.test.mjs} (\texttt{npm test}, wbudowany test runner Node.js, 8 przypadków) sprawdza dane wykresu ,,Procenty form w czasie'' dla ścieżki ankieta $\to$ 2 quizy $\to$ A $\to$ V $\to$ A: punkty tylko dla ankiety i policzonych quizów, zmiany stylu jako znaczniki (cofnięta oznaczona), etykiety osi ze zdarzeń i ,,bez zmian'' w wierszach zmian stylu. Kod klienta przechodzi lint (ESLint), kontrolę formatowania (Prettier) i build produkcyjny.

\subsection{Środowisko testowe}

\begin{center}
\begin{tabular}{ll}
\toprule
\textbf{Narzędzie} & \textbf{Wersja / uwagi} \\
\midrule
Python & 3.12+ \\
pytest & testy jednostkowe i integracyjne \\
FastAPI TestClient (httpx) & żądania HTTP do aplikacji \\
PostgreSQL & 17, baza \texttt{imsu\_test} w kontenerze Docker \\
Alembic & budowa schematu bazy testowej \\
Ruff & lint i formatowanie kodu serwera \\
Node.js, Playwright & przejście scenariusza pokazu (Chromium) \\
\bottomrule
\end{tabular}
\end{center}

\subsection{Przypadki testowe}

Poniższa tabela zawiera reprezentatywne przypadki testowe. Kolumna \emph{Typ}: \textbf{J} -- jednostkowy, \textbf{I} -- integracyjny (HTTP).

\begingroup
\scriptsize
\setlength{\tabcolsep}{2pt}
\begin{longtable}{p{1cm} >{\raggedright\arraybackslash}p{3.2cm} >{\raggedright\arraybackslash}p{1.6cm} p{0.6cm} >{\raggedright\arraybackslash}p{4.3cm} >{\raggedright\arraybackslash}p{4.3cm}}
\toprule
\textbf{ID} & \textbf{Nazwa} & \textbf{Moduł} & \textbf{Typ} & \textbf{Dane wejściowe / warunki} & \textbf{Oczekiwany wynik} \\
\midrule
\endfirsthead
\toprule
\textbf{ID} & \textbf{Nazwa} & \textbf{Moduł} & \textbf{Typ} & \textbf{Dane wejściowe / warunki} & \textbf{Oczekiwany wynik} \\
\midrule
\endhead
\midrule
\multicolumn{6}{r}{\footnotesize Kontynuacja na następnej stronie\ldots} \\
\endfoot
\bottomrule
\endlastfoot

TC-01 & Profil z ankiety (R1) & learning\_style & J &
  punkty A 11 / V 6 / K 3; V 9 / A 8 / K 3; V 20 / A 0 / K 0 &
  A 55 / V 30 / K 15; V 45 / A 40 / K 15; V 90 / A 5 / K 5 \faCheck \\
\addlinespace
TC-02 & Zmiana procentów (R2) & learning\_style & J &
  V 45 / A 40 / K 15, porażka w V; potem zaliczenie w A &
  V 35 / A 47 / K 18; V 28 / A 57 / K 15; suma 100, każda $\geq$ 5 \faCheck \\
\addlinespace
TC-03 & Tylko pierwsze podejście w formie & learning\_style & J, I &
  drugie podejście w tej samej formie i lekcji; podejście w innej formie &
  powtórka zapisana bez zmiany procentów; inna forma policzona \faCheck \\
\addlinespace
TC-04 & Sugestia innej formy (R5) & learning\_style & J, I &
  porażka bez potwierdzonych form; porażka przy potwierdzonej formie &
  druga forma ankiety (\texttt{survey}); forma potwierdzona (\texttt{confirmed}) z najlepszym zaliczeniem \faCheck \\
\addlinespace
TC-05 & Propozycja tylko dla formy potwierdzonej (R6) & learning\_style & J, I &
  wyższy procent formy bez zaliczenia; potem zaliczenie w tej formie &
  brak propozycji; propozycja z dowodem (lekcje, średnia, procent) \faCheck \\
\addlinespace
TC-06 & Odrzucona propozycja wraca (R6) & learning\_style & J, I &
  odrzucenie, potem 1 i 2 nowe policzone wyniki &
  po 1 wyniku brak propozycji, po 2 -- propozycja wraca \faCheck \\
\addlinespace
TC-07 & Granice okresu próbnego (R7) & learning\_style & J &
  różnica średnich $+10$, $+9$, $-9$, $-10$; brak średniej ,,przed'' &
  \texttt{helped}, \texttt{neutral}, \texttt{neutral}, \texttt{worse} + propozycja powrotu; \texttt{no\_baseline} \faCheck \\
\addlinespace
TC-08 & Ścieżka Kuby & learning\_style & I &
  porażka w V, zmiana formy, zaliczenie w A, przyjęcie propozycji &
  V 35 / A 47 / K 18 $\to$ V 28 / A 57 / K 15, propozycja V $\to$ A, okres próbny z ,,przed'' 40\% \faCheck \\
\addlinespace
TC-09 & Zmiana stylu czyści wybory formy & learning\_style & I &
  wybory formy w lekcji zaliczonej i niezaliczonej, zmiana stylu &
  wybór w niezaliczonej usunięty i opisany w zdarzeniu; w zaliczonej zostaje \faCheck \\
\addlinespace
TC-10 & Ankieta raz, remis & learning\_style & I &
  ponowne wysłanie ankiety; remis punktów &
  409 \texttt{SURVEY\_ALREADY\_DONE}; profil bez stylu, wybór przez \texttt{tie-choice} \faCheck \\
\addlinespace
\midrule
TC-11 & Nauczyciel czeka na zatwierdzenie & identity & I &
  rejestracja nauczyciela, endpoint nauczyciela &
  \texttt{approval\_status = pending}; 403 \texttt{TEACHER\_NOT\_APPROVED} \faCheck \\
\addlinespace
TC-12 & Zatwierdzanie tylko przez administratora & identity & I &
  nauczyciel lub uczeń wywołuje \texttt{POST /users/\{id\}/approve} &
  403; po decyzji administratora nauczyciel działa \faCheck \\
\addlinespace
TC-13 & Prośba o zapis & curriculum & I &
  prośba, powtórka, odrzucenie, ponowna prośba &
  \texttt{pending}; bez zmian; \texttt{rejected}; znów \texttt{pending} \faCheck \\
\addlinespace
TC-14 & Treść tylko dla przyjętych & curriculum & I &
  uczeń bez przyjętej prośby otwiera lekcję &
  403 \texttt{ENROLLMENT\_REQUIRED} \faCheck \\
\addlinespace
TC-15 & Blokada kolejności (R10) & curriculum & J, I &
  lekcja 2 przed zaliczeniem lekcji 1; zaliczenie; zmiana kolejności &
  403 \texttt{LESSON\_LOCKED} z wymaganą lekcją; odblokowanie; zaliczona zostaje dostępna \faCheck \\
\addlinespace
\midrule
TC-16 & Gotowość lekcji (R11) & authoring & J, I &
  nowa lekcja; pula 3 pytań; publikacja niekompletnej lekcji &
  lista braków z kodami; \texttt{QUIZ\_TOO\_FEW\_QUESTIONS} (3, min. 5); 409 \texttt{LESSON\_NOT\_READY} \faCheck \\
\addlinespace
TC-17 & Opublikowana lekcja zostaje kompletna & content & I &
  usunięcie karty PDF z wersji K opublikowanej lekcji &
  409 \texttt{PUBLISHED\_LESSON\_INCOMPLETE} z listą braków \faCheck \\
\addlinespace
TC-18 & Plik innego nauczyciela & content & I &
  \texttt{PUT /lessons/\{id\}/versions/V}: załącznik z cudzym plikiem &
  400 \texttt{FILE\_NOT\_FOUND} \faCheck \\
\addlinespace
\midrule
TC-19 & Losowanie z puli (R9) & assessment & J, I &
  pula 8, pierwsze podejście i powtórka; odświeżenie strony &
  5 pytań bez klucza; powtórka: 3 nowe + 2 stare; ten sam otwarty zestaw \faCheck \\
\addlinespace
TC-20 & Odpowiedzi spoza zestawu & assessment & I &
  odpowiedź na pytanie spoza wylosowanego zestawu; drugie wysłanie &
  400 \texttt{QUIZ\_ANSWER\_MISMATCH}; 409 \texttt{DRAW\_NOT\_OPEN} \faCheck \\
\addlinespace
TC-21 & Wyniki z wyjaśnieniami & assessment & I &
  wysłanie zestawu &
  wynik, poprawne odpowiedzi i wyjaśnienie przy każdym pytaniu, \texttt{QuizFeedback} \faCheck \\
\addlinespace
\midrule
TC-22 & Statystyki i szczegóły ucznia & learning\_style & I &
  autor i inny nauczyciel pobierają statystyki kursu i szczegóły ucznia &
  autor: skuteczność form, rozkład stylów; inny nauczyciel -- 404, uczeń -- 403 \faCheck \\
\addlinespace
TC-23 & Usunięte endpointy v2 & learning\_style & I &
  \texttt{personalized-blocks}, \texttt{blocks/\{id\}/open}, \texttt{proposal/accept} &
  404 \faCheck \\
\addlinespace
TC-24 & Idempotentny seed i reset & demo & I &
  dwukrotne uruchomienie seeda; pokaz, potem \texttt{-{}-reset-students} &
  drugie uruchomienie nic nie zmienia; stany kont wracają do startowych \faCheck \\
\addlinespace
\midrule
TC-25 & Zmiana właściciela kursu & curriculum & I &
  admin: \texttt{PUT /courses/\{id\}/owner} na zatwierdzonego nauczyciela; ten sam właściciel; nauczyciel \texttt{pending} / \texttt{rejected}; uczeń, admin, brak konta; wywołanie przez nauczyciela &
  nowy autor w liście i katalogu, poprzedni -- 404, zapisy i stany lekcji zostają, wiersz w \texttt{course\_owner\_changes}; bez zmian i bez wpisu; 409 \texttt{TEACHER\_NOT\_APPROVED}; 404 \texttt{USER\_NOT\_FOUND}; 403 \faCheck \\
\addlinespace
TC-26 & Pliki po zmianie właściciela & content & I &
  nowy autor zapisuje wersję z plikiem poprzedniego autora użytym w lekcji (także w innej wersji); plik poprzedniego autora nieużyty w lekcji &
  zapis przechodzi; 400 \texttt{FILE\_NOT\_FOUND} \faCheck \\
\addlinespace
TC-27 & Cofnięcie zmiany stylu (R7a) & learning\_style & J, I &
  styl A (2 lekcje po 100\%) $\to$ V $\to$ A bez lekcji; A $\to$ V $\to$ K bez lekcji; zmiana po policzonej lekcji; cofnięcie po przyjętej propozycji &
  okres \texttt{cancelled}, brak okresu ,,V $\to$ A'', zdarzenie \texttt{style\_change\_reverted}, profil bez okresu; okres A $\to$ K z ,,przed'' z A i \texttt{replaced\_trial\_id}; nowy okres jak w R7 \faCheck \\
\addlinespace
TC-28 & Stany lekcji ucznia & learning\_style & I &
  \texttt{GET /courses/\{id\}/study-states}: nowy uczeń; sam wybór formy; wynik niezaliczony; lekcja zaliczona; nieopublikowana lekcja; nauczyciel i uczeń nieprzyjęty &
  \texttt{started = false} (,,Rozpocznij''); \texttt{started}, \texttt{chosen\_form}, \texttt{last\_percent} (,,Kontynuuj''); \texttt{passed}, \texttt{best\_percent} (,,Powtórz''); tylko opublikowane; 403 \faCheck \\
\addlinespace
TC-29 & Forma równa stylowi & learning\_style & I &
  wybór formy równej stylowi; wybór w zaliczonej lekcji, potem zmiana stylu na tę formę &
  wybór usunięty; \texttt{form\_source = style} (,,Twój styl''), nie \texttt{choice} \faCheck \\
\addlinespace
TC-30 & \texttt{null} w \texttt{PATCH} kursu i lekcji & curriculum & I &
  \texttt{null} we wstępie, podsumowaniu, zarysie, wnioskach szkicu; \texttt{null} w tytule lub opisie; pole pominięte &
  pole wyczyszczone; 422 \texttt{VALIDATION\_ERROR}; pole bez zmian \faCheck \\
\addlinespace
TC-31 & Zestaw pytań a forma lekcji & assessment & I &
  otwarty zestaw w formie V, zmiana formy lekcji na A; powrót do tej samej formy; zestaw sprzed migracji (\texttt{form = null}) &
  stary zestaw usunięty, nowy z \texttt{form = A}; bez zmian; zestaw zostaje \faCheck \\
\addlinespace
TC-32 & Historia ze zdarzeniami, pochwała w dzienniku & learning\_style & I &
  ankieta, policzone i niepoliczone quizy, zmiana ręczna, cofnięcie, seria 2 zaliczeń w jednej formie &
  punkty historii z \texttt{event} (bez punktów dla niepoliczonych wyników i pochwały), numer lekcji; zdarzenie \texttt{praise} \faCheck \\
\addlinespace
TC-33 & Nauczyciele kursów i nagrania demo & demo & I &
  nowa baza; \texttt{-{}-reset-students} po przepięciu Angielskiego na Piotra; \texttt{-{}-refresh-demo-media} &
  każdy kurs u swojego nauczyciela, bez wpisów zmian właściciela; właściciele przywróceni (2 wpisy); wersje A z 3--4 nagraniami z plikiem i transkrypcją, ostatnie ,,Powtórka'' \faCheck \\

\end{longtable}
\endgroup

\subsection{Scenariusze testowe}

\subsubsection{ST-01: Od ankiety do sugestii}
\textbf{Cel:} sprawdzenie ścieżki ucznia od ankiety do pierwszej reakcji na porażkę.
\begin{enumerate}[leftmargin=*]
    \item Uczeń bez ankiety otwiera lekcję -- 403 \texttt{SURVEY\_REQUIRED}.
    \item Uczeń wypełnia ankietę -- powstaje profil procentowy i styl (TC-01); druga ankieta jest odrzucona (TC-10).
    \item Uczeń otwiera lekcję -- dostaje wersję w formie swojego stylu; kolejna lekcja jest zablokowana (TC-15).
    \item Uczeń rozwiązuje quiz z puli (TC-19) i nie zalicza -- procent formy spada o 10 (TC-02), a wynik zawiera wyjaśnienia (TC-21) i sugestię innej formy (TC-04).
    \item Uczeń wybiera ,,Przerób tę lekcję w formie X'' -- dostaje wersję X i nowy zestaw pytań; powtórka w tej samej formie nie zmienia procentów (TC-03).
\end{enumerate}

\subsubsection{ST-02: Propozycja zmiany stylu i okres próbny}
\textbf{Cel:} sprawdzenie reakcji systemu na wyniki w innej formie (ścieżka Kuby, TC-08).
\begin{enumerate}[leftmargin=*]
    \item Uczeń zalicza lekcję w formie innej niż styl -- forma jest potwierdzona, ma wyższy procent, powstaje propozycja z dowodem (TC-05).
    \item Uczeń odrzuca propozycję -- wraca ona po 2 nowych policzonych wynikach (TC-06).
    \item Uczeń przyjmuje propozycję -- styl się zmienia, wybory formy w niezaliczonych lekcjach są czyszczone (TC-09), zaczyna się okres próbny.
    \item Po 2 policzonych wynikach w nowej formie okres kończy się wynikiem zależnym od różnicy średnich; spadek o 10 lub więcej daje propozycję powrotu (TC-07).
\end{enumerate}

\subsubsection{ST-03: Zapisy na kurs i zatwierdzanie nauczycieli}
\textbf{Cel:} sprawdzenie dostępu do kursów i kont nauczycieli.
\begin{enumerate}[leftmargin=*]
    \item Nauczyciel rejestruje się -- konto czeka na zatwierdzenie, endpointy nauczyciela zwracają 403 (TC-11); administrator zatwierdza konto (TC-12).
    \item Nauczyciel publikuje lekcję dopiero, gdy jest kompletna (TC-16, TC-17), a potem kurs.
    \item Uczeń prosi o zapis -- do decyzji nie ma dostępu do treści (TC-13, TC-14).
    \item Autor przyjmuje prośbę -- uczeń widzi stany lekcji; po wycofaniu publikacji kursu przyjęty uczeń zachowuje dostęp.
    \item Administrator przekazuje kurs innemu zatwierdzonemu nauczycielowi -- prośby, zapisy i wyniki przechodzą razem z kursem, poprzedni autor traci dostęp, nowy może zapisywać lekcje z plikami poprzedniego (TC-25, TC-26).
\end{enumerate}

\subsubsection{ST-04: Zmiana stylu cofnięta przed pierwszą lekcją}
\textbf{Cel:} sprawdzenie błędu zgłoszonego po teście 12.11 (TC-27, TC-32) -- dokładnie ścieżka z testu.
\begin{enumerate}[leftmargin=*]
    \item Uczeń ze stylem A zalicza 2 lekcje w A po 100\%.
    \item Uczeń zmienia styl ręcznie na V -- zaczyna się okres próbny ze średnią ,,przed'' 100\% w formie A.
    \item Bez żadnej lekcji uczeń wraca do A -- okres próbny dostaje \texttt{cancelled}, nowy okres ,,V $\to$ A'' bez danych do porównania nie powstaje, dziennik ma wpis o cofnięciu.
    \item ,,Moje postępy'': wykres ma 3 punkty (ankieta i 2 quizy) oraz 2 znaczniki zmian stylu, z których drugi jest przekreślony (,,cofnięta''); tabela ma kolumnę ,,Zdarzenie'' i ,,bez zmian'' w wierszach zmian stylu (test komponentu \texttt{chart-data}).
    \item Wariant: zamiast powrotu do A -- przejście na K; okres A $\to$ K zachowuje średnią ,,przed'' z formy A.
\end{enumerate}

\subsection{Podsumowanie pokrycia testami}

\begin{center}
\begin{tabular}{lrrr}
\toprule
\textbf{Moduł} & \textbf{Jednostkowe} & \textbf{Integracyjne} & \textbf{Łącznie} \\
\midrule
\texttt{identity}        & 6   & 73  & 79 \\
\texttt{curriculum}      & 21  & 127 & 148 \\
\texttt{content}         & 35  & 45  & 80 \\
\texttt{learning\_style} & 81  & 65  & 146 \\
\texttt{assessment}      & 14  & 23  & 37 \\
\texttt{authoring}       & --  & 16  & 16 \\
\texttt{demo}            & --  & 19  & 19 \\
\texttt{core} (TipTap, błędy, stan) & 10 & 5 & 15 \\
\midrule
\textbf{Razem}           & \textbf{167} & \textbf{373} & \textbf{540} \\
\bottomrule
\end{tabular}
\end{center}

Tabela podaje liczbę przypadków wykonywanych przez pytest (z parametryzacją; 386 funkcji testowych), policzoną poleceniem \texttt{pytest -{}-collect-only -q}. W porównaniu z wersją sprzed poprawek z 12.11 (473) doszły m.in.: \nolinkurl{tests/curriculum/test_owner.py} (17, zmiana właściciela), przypadki R7a w \nolinkurl{test_rules.py} i \nolinkurl{test_adaptation_flow.py}, stany lekcji i historia ze zdarzeniami w \nolinkurl{test_insights_api.py}, forma równa stylowi w \nolinkurl{test_study_api.py}, \texttt{null} w \texttt{PATCH} kursu i lekcji, zestaw pytań z formą w \nolinkurl{test_quiz_api.py}, pliki po zmianie właściciela i nagranie z transkrypcją w \nolinkurl{test_versions.py}, kod w dokumentach TipTap (\nolinkurl{tests/test_rich_text.py}, 10) oraz nauczyciele kursów, nagrania i \texttt{-{}-refresh-demo-media} w teście seeda. Sam moduł reguł \nolinkurl{learning_style/rules.py} ma \textbf{73} przypadki testowe w \nolinkurl{tests/learning_style/test_rules.py}. Test seeda sprawdza stany wszystkich kont demonstracyjnych (procenty, propozycje, okres próbny, prośbę \texttt{pending}, nauczyciela \texttt{pending}), właścicieli kursów, wersje A z nagraniami i transkrypcjami, bloki kodu, idempotencję, \texttt{-{}-reset-students} i \texttt{-{}-refresh-demo-media}. Przejście Playwright \nolinkurl{frontend/e2e/demo-tour.mjs} obejmuje wszystkie konta demonstracyjne w kolejności scenariusza pokazu (50 zrzutów), a test komponentu \texttt{chart-data} -- 8 przypadków.

\section{Instrukcja uruchomienia aplikacji lokalnie}
\label{sec:setup}

Baza danych działa w kontenerze Docker, a serwer i klient są uruchamiane lokalnie. Polecenia podano dla systemu Windows (wiersz poleceń \texttt{cmd}); każde polecenie należy wykonać osobno. Instrukcja opisuje uruchomienie od zera na nowej maszynie.

\subsection{Wymagania}

\begin{itemize}[leftmargin=*]
    \item \textbf{Docker Desktop} -- baza danych PostgreSQL (\url{https://www.docker.com/products/docker-desktop/})
    \item \textbf{Python 3.12} lub nowszy -- serwer (\url{https://www.python.org/downloads/})
    \item \textbf{Node.js 20} lub nowszy wraz z \textbf{npm} -- klient (\url{https://nodejs.org/})
    \item \textbf{Git} -- pobranie repozytorium (\url{https://git-scm.com/downloads})
\end{itemize}

Sprawdzenie zainstalowanych narzędzi:

\begin{lstlisting}
docker --version
python --version
node --version
git --version
\end{lstlisting}

\subsection{Krok 1: Pobranie repozytorium}

\begin{lstlisting}
git clone https://github.com/PatrykKul/aplikacjaIMSU.git
cd aplikacjaIMSU
\end{lstlisting}

\subsection{Krok 2: Zmienne środowiskowe}

Projekt korzysta z trzech plików ze zmiennymi środowiskowymi. Każdy powstaje przez skopiowanie dołączonego wzoru (pliki \texttt{.env} nie trafiają do repozytorium):

\begin{lstlisting}
copy .env.example .env
copy backend\.env.example backend\.env
copy frontend\.env.example frontend\.env.local
\end{lstlisting}

\textbf{Baza danych} (\texttt{.env} w katalogu głównym): \texttt{POSTGRES\_USER}, \texttt{POSTGRES\_PASSWORD}, \texttt{POSTGRES\_DB}, \texttt{POSTGRES\_PORT} (domyślnie 5433).

\textbf{Serwer} (\texttt{backend\textbackslash.env}): \texttt{DATABASE\_URL}, \texttt{TEST\_DATABASE\_URL} (baza \texttt{imsu\_test}), \texttt{CORS\_ORIGINS} (adres klienta \url{http://localhost:3001}), \texttt{JWT\_SECRET}, \texttt{UPLOAD\_DIR} (katalog wgranych plików, względem \texttt{backend}).

\textbf{Klient} (\texttt{frontend\textbackslash.env.local}): \texttt{NEXT\_PUBLIC\_API\_URL} (adres serwera \url{http://localhost:8001}).

\begin{tcolorbox}[colback=yellow!10!white,colframe=orange!75!black,title=\faExclamationTriangle\ Klucz JWT]
\texttt{JWT\_SECRET} musi mieć co najmniej 32 znaki -- bez niego serwer się nie uruchomi. Własny losowy klucz można wygenerować poleceniem \verb|python -c "import secrets; print(secrets.token_hex(32))"|.
\end{tcolorbox}

\subsection{Krok 3: Baza danych}

Uruchom Docker Desktop, a następnie w katalogu głównym repozytorium:

\begin{lstlisting}
docker compose up -d
\end{lstlisting}

Polecenie uruchamia PostgreSQL 17 na porcie \textbf{5433}. Przy pierwszym starcie kontener tworzy dwie bazy: \texttt{imsu} (aplikacja) i \texttt{imsu\_test} (testy).

\subsection{Krok 4: Serwer}

Instalacja zależności w środowisku wirtualnym (tylko przy pierwszym uruchomieniu):

\begin{lstlisting}
cd backend
python -m venv venv
venv\Scripts\python.exe -m pip install -r requirements.txt
\end{lstlisting}

Utworzenie tabel i wstawienie 20 pytań ankiety (migracje), a następnie wczytanie danych demonstracyjnych:

\begin{lstlisting}
venv\Scripts\python.exe -m alembic upgrade head
venv\Scripts\python.exe -m app.demo.seed
\end{lstlisting}

Seed tworzy konta, 3 kursy po 4 opublikowane lekcje -- każdy kurs u innego nauczyciela, który sam wgrywa jego pliki (wersje V / A / K z obrazami, nagraniami MP3 i kartami PDF, razem 79 plików; pule po 8 pytań) -- oraz historie uczniów. Wersja słuchowa każdej lekcji ma krótki opis, 3 nagrania po 30--50 s zatytułowane jak sekcje wersji V i nagranie ,,Powtórka'' (20--30 s), każde z transkrypcją; kod w lekcjach Programowania (i rachunki w Matematyce) jest zapisany jako bloki kodu. Historie powstają przez prawdziwe serwisy (ankieta, losowanie, wysłanie odpowiedzi, decyzje), więc procenty i dziennik są prawdziwe; seed tylko przesuwa daty zdarzeń o dni wstecz. Seed jest \textbf{idempotentny}: drugie uruchomienie niczego nie zmienia i wypisuje ,,bez zmian''. Na końcu wypisuje ,,Stan kont'' i ,,Skuteczność form'', np.:

\begin{quote}\ttfamily\small
uczen.demo@imsu.pl: styl V, V 45 / A 40 / K 15, Matematyka 2/4\\
Matematyka: funkcja liniowa (Anna Nowak): V 60\% (3 ucz.) \textperiodcentered{} A 86\% (3 ucz.) \textperiodcentered{} K 70\% (2 ucz.)
\end{quote}

Baza utworzona przed poprawkami z 12.11 (jedno nagranie w wersji A, wszystkie kursy u jednej nauczycielki): po \texttt{alembic upgrade head} zwykłe uruchomienie seeda przepina Angielski i Programowanie na ich nauczycieli (wypisuje ,,zmieniono autora kursu \ldots'' i zapisuje zmianę w \texttt{course\_owner\_changes}), a jednorazowo:

\begin{lstlisting}
venv\Scripts\python.exe -m app.demo.seed --refresh-demo-media
\end{lstlisting}

zapisuje od nowa treść kursów demo z \nolinkurl{app/demo/courses.py} (wersje A z nowymi nagraniami i transkrypcjami, bloki kodu, podsumowania kursów). Obrazy V i karty K zostają; zmiany nauczyciela w treści kursów demo są nadpisywane.

Uruchomienie serwera:

\begin{lstlisting}
venv\Scripts\python.exe -m uvicorn app.main:app --reload --reload-dir app --port 8001
\end{lstlisting}

Serwer działa pod adresem \url{http://localhost:8001}, a dokumentacja API (Swagger) pod \url{http://localhost:8001/docs}.

\subsection{Krok 5: Klient}

W drugim oknie wiersza poleceń, z katalogu głównego repozytorium (\texttt{npm ci} instaluje dokładnie wersje z \texttt{package-lock.json}):

\begin{lstlisting}
cd frontend
npm ci
npm run dev
\end{lstlisting}

Aplikacja jest dostępna pod adresem \url{http://localhost:3001}. Wersję produkcyjną klienta uruchamia się poleceniami \texttt{npm run build} i \texttt{npm run start} (ten sam port).

\subsection{Krok 6: Weryfikacja}

\begin{enumerate}[leftmargin=*]
    \item Adres \url{http://localhost:8001/health} zwraca odpowiedź:\\ \verb|{"status": "ok", "database": "ok"}|
    \item Pod adresem \url{http://localhost:3001} widać stronę startową z przyciskami ,,Zaloguj się'' i ,,Załóż konto''
    \item Logowanie na konto demonstracyjne działa (hasło wszystkich kont: \texttt{Demo123!}):
\end{enumerate}

\begingroup
\small
\begin{longtable}{|p{6.4cm}|p{2.4cm}|p{5.7cm}|}
\hline
\textbf{Konto} & \textbf{Rola} & \textbf{Stan startowy} \\ \hline
\endhead
\texttt{admin.demo@imsu.pl} & administrator & zatwierdza nauczyciela, zarządza kontami, zmienia właściciela kursu \\ \hline
\texttt{nauczyciel.demo@imsu.pl} & nauczyciel & Anna Nowak, autorka Matematyki: 7 uczniów i 1 prośba o zapis \\ \hline
\texttt{nauczyciel.angielski@imsu.pl} & nauczyciel & Katarzyna Wróbel, autorka Angielskiego (Antoni, Szymon) \\ \hline
\texttt{nauczyciel.programowanie@imsu.pl} & nauczyciel & Paweł Sikora, autor Programowania (Marta, Filip) \\ \hline
\texttt{nauczyciel.nowy@imsu.pl} & nauczyciel & Piotr Zieliński, konto czeka na zatwierdzenie \\ \hline
\texttt{uczen.nowy@imsu.pl} & uczeń & Ola Kowalska, bez ankiety (wymuszenie ankiety) \\ \hline
\texttt{uczen.zapis@imsu.pl} & uczeń & Zosia Mazur, A 50 / V 30 / K 20, prośba o zapis do Matematyki \\ \hline
\texttt{uczen.demo@imsu.pl} & uczeń & Kuba Wiśniewski, V 45 / A 40 / K 15, Matematyka 2/4 -- ,,tuż przed sugestią'' \\ \hline
\texttt{uczen.zmiana@imsu.pl} & uczeń & Marta Lewandowska, po zmianie K $\to$ V: ,,Zmiana pomogła: średnia 40\% $\to$ 90\%'' \\ \hline
\texttt{uczen1@imsu.pl} \ldots\ \texttt{uczen6@imsu.pl} & uczniowie & wyniki w Matematyce (skuteczność form w panelu kursu); jedna otwarta propozycja; \texttt{uczen3} (Szymon) -- styl A do pokazu cofnięcia zmiany stylu (R7a) \\ \hline
\caption{Konta demonstracyjne}
\end{longtable}
\endgroup

Stan startowy wszystkich kont demonstracyjnych (także po pokazie) przywraca polecenie (z katalogu \texttt{backend}); kursy i ich treść zostają, a kurs demo przepięty na pokazie na inne konto demo wraca do swojego nauczyciela (przez tę samą funkcję co w aplikacji, z wpisem w \texttt{course\_owner\_changes}):
\begin{lstlisting}
venv\Scripts\python.exe -m app.demo.seed --reset-students
\end{lstlisting}

Własne konto administratora tworzy się w konsoli (hasło jest podawane interaktywnie):
\begin{lstlisting}
venv\Scripts\python.exe -m app.modules.identity.cli create-admin --email admin@example.com --name "Administrator"
\end{lstlisting}

\subsection{Uruchomienie testów i kontrola kodu}

Serwer (katalog \texttt{backend}, baza danych musi działać):
\begin{lstlisting}
venv\Scripts\python.exe -m pytest
venv\Scripts\python.exe -m ruff check .
venv\Scripts\python.exe -m ruff format --check .
\end{lstlisting}

Klient (katalog \texttt{frontend}):
\begin{lstlisting}
npm run lint
npm run format:check
npm test
npm run build
\end{lstlisting}
\texttt{npm run build} nie należy uruchamiać przy działającym \texttt{npm run dev} na tym samym katalogu \texttt{.next} -- najpierw zatrzymać serwer deweloperski.

Przejście scenariusza pokazu w przeglądarce (serwer i klient uruchomione, świeże dane demo). Playwright nie jest zależnością projektu -- instaluje się go w osobnym katalogu i wskazuje zmienną \texttt{PLAYWRIGHT\_MODULE}:
\begin{lstlisting}
npm i playwright
npx playwright install chromium
set PLAYWRIGHT_MODULE=<katalog>\node_modules\playwright\index.js
node e2e/demo-tour.mjs
\end{lstlisting}
Skrypt zapisuje 50 zrzutów ekranów do \nolinkurl{docs/zrzuty-v3} i na końcu wypisuje listę problemów (kod wyjścia 1, jeśli jakiś wystąpił). Po przejściu należy wykonać \texttt{-{}-reset-students}, bo skrypt zmienia stan kont.

\subsection{Rozwiązywanie problemów}

\subsubsection{Endpoint \texttt{/health} zwraca błąd bazy danych}
Baza danych nie działa. Uruchom Docker Desktop i ponownie wykonaj \texttt{docker compose up -d} w katalogu głównym repozytorium (konieczne także po restarcie komputera).

\subsubsection{Serwer nie startuje: błąd \texttt{jwt\_secret}}
W pliku \texttt{backend\textbackslash.env} brakuje zmiennej \texttt{JWT\_SECRET} lub ma ona mniej niż 32 znaki.

\subsubsection{Port 3001, 8001 lub 5433 jest zajęty}
Błąd \texttt{EADDRINUSE} przy \texttt{npm run dev} oznacza zwykle, że klient już działa w innym oknie. Port bazy można zmienić zmienną \texttt{POSTGRES\_PORT} (i odpowiednio adresy w \texttt{backend\textbackslash.env}).

\subsubsection{Brak pytań ankiety lub kursów demonstracyjnych}
Wykonaj migracje (\texttt{alembic upgrade head}), a potem seed (\texttt{app.demo.seed}). Seed można uruchamiać wielokrotnie -- dokłada tylko brakujące dane.

\subsubsection{Baza z wersji 2}
Migracja v3 nie przenosi danych uczenia z v2 (profile, próby quizów, quizy); konta, kursy, lekcje i pliki zostają, a lekcje stają się nieopublikowane. Najprościej przebudować lokalną bazę: \texttt{alembic downgrade base}, \texttt{alembic upgrade head}, seed.

\subsubsection{Wersja A ma jedno nagranie albo kursy nie mają swoich nauczycieli}
Baza pochodzi sprzed poprawek z 12.11. Wykonaj \texttt{alembic upgrade head}, seed (przepina kursy na nauczycieli) i raz \texttt{app.demo.seed -{}-refresh-demo-media} (nowe nagrania z transkrypcjami, bloki kodu).

\subsubsection{Seed utworzył drugi kurs o tym samym tytule}
Kurs demo został przepięty na konto spoza kont demonstracyjnych (np. nauczyciela zarejestrowanego na pokazie), więc seed go nie rozpoznał. Przed resetem przepnij kurs z powrotem (,,Zmień właściciela'') albo zarchiwizuj duplikat.

\subsubsection{Polecenia z \texttt{venv\textbackslash Scripts} nie działają po przeniesieniu projektu}
Pliki \texttt{.exe} w \texttt{venv\textbackslash Scripts} mają zapisane bezwzględne ścieżki. Uruchamiaj narzędzia przez \texttt{python -m ...} (jak w tej instrukcji) albo utwórz środowisko od nowa.

\subsection{Struktura projektu}

\begin{lstlisting}
aplikacjaIMSU/
  docker-compose.yml        # baza danych PostgreSQL
  docker/                   # skrypt tworzacy baze testowa
  backend/                  # serwer - FastAPI
    requirements.txt        # zaleznosci Python
    alembic/                # migracje bazy danych
    app/
      main.py               # punkt wejscia aplikacji
      core/                 # konfiguracja, sesja bazy, bledy
      modules/
        identity/           # konta, role, logowanie, zatwierdzanie
        curriculum/         # kursy, lekcje, zapisy, blokada (R10)
        content/            # wersje V/A/K lekcji, pliki
        learning_style/     # ankieta, profil, reguly R1-R8, dziennik
        assessment/         # pula pytan, losowanie (R9), ocenianie
        authoring/          # gotowosc (R11) i publikacja lekcji
      demo/                 # dane demonstracyjne (seed, kursy, zasoby)
    tests/                  # testy pytest
  frontend/                 # klient - Next.js
    e2e/demo-tour.mjs       # przejscie scenariusza pokazu (Playwright)
    src/
      app/                  # trasy (strony)
      features/             # auth, courses, content, learning-style,
                            #   assessment, progress, teach, users...
      shared/               # komponenty UI, klient API, uklad strony
  docs/                     # dokumentacja projektu
\end{lstlisting}

\subsection{Przydatne polecenia}

\begin{lstlisting}
# Baza danych (katalog glowny)
docker compose up -d                                   # start bazy
docker compose down                                    # zatrzymanie bazy

# Serwer (katalog backend)
venv\Scripts\python.exe -m alembic upgrade head        # migracje
venv\Scripts\python.exe -m app.demo.seed               # dane demo
venv\Scripts\python.exe -m app.demo.seed --reset-students  # reset kont
venv\Scripts\python.exe -m app.demo.seed --refresh-demo-media  # nowa tresc demo
venv\Scripts\python.exe -m uvicorn app.main:app --port 8001   # start
venv\Scripts\python.exe -m pytest                      # testy

# Klient (katalog frontend)
npm ci                                                 # zaleznosci
npm run dev                                            # start (deweloperski)
npm run build                                          # build produkcyjny
npm run start                                          # start (produkcyjny)
\end{lstlisting}


\end{document}